% PRD-style article: Quasi-local Observable Domain Reconstruction in S^3, v29 final submission package
\documentclass[aps,prd,onecolumn,superscriptaddress,longbibliography,nofootinbib,notitlepage]{revtex4-2}

\usepackage{amsmath,amssymb,amsfonts}
\usepackage{mathtools}
\usepackage{physics}
\usepackage{graphicx}
\usepackage{xcolor}
\usepackage{hyperref}
\usepackage{bm}
\usepackage{booktabs}
\usepackage{array}
\usepackage{microtype}
\emergencystretch=2em

% ----- Single-column journal layout layer -----
\usepackage[a4paper,top=24mm,bottom=25mm,left=22mm,right=22mm]{geometry}
\usepackage{etoolbox}
\usepackage{colortbl}
\usepackage[most]{tcolorbox}
\usepackage{fancyhdr}
\usepackage{lastpage}
\usepackage{enumitem}
\usepackage{needspace}
\usepackage{titlesec}
\definecolor{JournalBlue}{HTML}{143D59}
\definecolor{JournalGold}{HTML}{B58B2B}
\definecolor{JournalInk}{HTML}{202124}
\definecolor{JournalSoft}{HTML}{F6F8FB}
\definecolor{JournalRule}{HTML}{D9E2EC}
\definecolor{JournalBox}{HTML}{EEF4FA}
\color{JournalInk}
\hypersetup{linkcolor=JournalBlue,citecolor=JournalBlue,urlcolor=JournalBlue}
\setlength{\parindent}{1.1em}
\setlength{\parskip}{0.16em}
\setlist{itemsep=0.15em,topsep=0.2em}
\renewcommand{\arraystretch}{1.17}
\newcommand{\journalrule}{\par\vspace{0.45em}{\color{JournalGold}\hrule height 0.8pt}\vspace{0.65em}}
\titleformat{\section}{\needspace{7\baselineskip}\large\bfseries\color{JournalBlue}}{\thesection}{0.75em}{}[\vspace{-0.25em}{\color{JournalGold}\titlerule[0.6pt]}]
\titleformat{\subsection}{\needspace{5\baselineskip}\normalsize\bfseries\color{JournalBlue}}{\thesubsection}{0.65em}{}
\titleformat{\subsubsection}{\needspace{4\baselineskip}\normalsize\bfseries\itshape\color{JournalBlue}}{\thesubsubsection}{0.55em}{}
\titlespacing*{\section}{0pt}{1.35em}{0.65em}
\titlespacing*{\subsection}{0pt}{1.0em}{0.4em}
\tcbset{colback=JournalSoft,colframe=JournalBlue!55!black,boxrule=0.55pt,arc=1.4mm,left=1.1mm,right=1.1mm,top=0.8mm,bottom=0.8mm}
\AtBeginEnvironment{quote}{\begin{tcolorbox}[enhanced,breakable,colback=JournalBox,colframe=JournalBlue!45!black,borderline west={1.5pt}{0pt}{JournalGold}]}
\AtEndEnvironment{quote}{\end{tcolorbox}}
\makeatletter
\pretocmd{\table}{\small\rowcolors{2}{JournalSoft}{white}}{}{}
\apptocmd{\endtable}{\normalsize}{}{}
\makeatother
\setlength{\headheight}{15pt}
\pagestyle{fancy}
\fancyhf{}
\fancyhead[L]{\footnotesize\textsc{Conditional Reconstruction within Compact $S^3$ Geometry}}
\fancyhead[R]{\footnotesize Preece \& Batenin}
\fancyfoot[C]{\footnotesize\thepage\,/\,\pageref{LastPage}}
\renewcommand{\headrulewidth}{0.4pt}
\renewcommand{\footrulewidth}{0pt}


\hypersetup{
  colorlinks=true,
  linkcolor=blue!70!black,
  citecolor=blue!70!black,
  urlcolor=blue!70!black,
  pdftitle={Conditional Reconstruction of a Quasi-Local Observable Domain within Compact S3 Geometry},
  pdfauthor={Andrei Preece and Boris Batenin},
  pdfsubject={Conditional compact S3 cosmological reconstruction and invariant diagnostic registry},
  pdfkeywords={S3 geometry, compact cosmology, FLRW curvature, hyperspherical harmonics, CMB, topology}
}

% custom commands
\newcommand{\SR}{S^3_R}
\newcommand{\RS}{R_{\mathcal{S}^3}}
\newcommand{\Gly}{\,\mathrm{Gly}}
\newcommand{\Gyr}{\,\mathrm{Gyr}}
\newcommand{\Mly}{\,\mathrm{Mly}}
\newcommand{\phivarphi}{\varphi}
\newcommand{\GammaT}{\Gamma(T)}

\begin{document}
\raggedbottom

\title{Conditional Reconstruction of a Quasi-Local Observable Domain\\ within Compact \texorpdfstring{$S^3$}{S3} Geometry}

\author{Andrei Preece}
\email{andre.preece@web.de}
\homepage[ORCID: ]{https://orcid.org/0009-0004-5978-795X}
\affiliation{Independent Researcher, Berlin, Germany}

\author{Boris Batenin}
\email{batenin@web.de}
\homepage[ORCID: ]{https://orcid.org/0009-0000-7027-8583}
\affiliation{Independent Researcher, Berlin, Germany}

\date{June 9, 2026}

\begin{abstract}
We present a conditional reconstruction of a quasi-local observable domain inside a
compact, simply connected, positively curved round three-sphere $S^3_R$. The
reconstruction is not a topology detection, dark-sector derivation, or residual-source
completion. It is a background-reconstructive and benchmark-ready compact-geometry
framework. In standard FLRW convention, closed positive spatial curvature corresponds to
$\Omega_K<0$. We use the fiducial modulus $|\Omega_K|=0.0007$ only as a curvature-radius
benchmark and set $\Omega_K^{S^3}=-0.0007$ for the closed branch. With
$H_0=67.4\,\mathrm{km\,s}^{-1}\mathrm{Mpc}^{-1}$ this gives
$R_{\mathcal S^3}\simeq548.3\,\mathrm{Gly}$; machine-precision benchmark values are
kept only in the numerical table.

The exact geometry then supplies $D_c=\pi R$, $C=2\pi R$,
$V_{S^3}=2\pi^2R^3$, $K=R^{-2}$, and ${}^{(3)}R=6R^{-2}$. The observable particle-horizon
scale used in the benchmark, $D_{\rm particle}\simeq46.13\,\mathrm{Gly}$, is treated as
background dependent rather than as a theorem of $S^3$ geometry. Its geodesic-ball
fraction is $\Gamma_{\rm obs}=1.2614\times10^{-4}\simeq0.0126\%$. The static-carrier
comparison $ct_0/R\simeq1.44^\circ$ is only a diagnostic scale; FLRW causal propagation is
governed by conformal distance. The finite radius is therefore a curvature scale, not a
boundary: $\partial S^3_R=\varnothing$, while observer-defined subdomains can have
relational reconstruction boundaries.

The predictive content is registered through the chain
$|\Omega_K|\to R_c\to\{D_M,H,dV/dz\,d\Omega\}\to\{\beta,\Delta_{\ell\beta}^X\}\to
\{C_\ell^{TT,TE,EE},P(k,z),\xi\}\to\Delta\ln\mathcal L$ or $\Delta\ln Z$. Round
$S^3_R$ has no invariant spatial direction under the full $SO(4)$ symmetry; any axis,
ring normal, or preferred frame must enter through data, fields, survey windows, initial
conditions, or quotient topology $S^3/\Gamma$. Thus the paper's empirical target is not a
preferred direction but the replacement of the flat continuum by the $S^3$ eigenvalue
ladder, degeneracy law, addition kernel, windowed responses, and covariance matrices.
All amplitude claims remain transfer- and likelihood-dependent.
\end{abstract}

\maketitle
\journalrule

\section{Introduction and Purpose}
\label{sec:intro}

The aim of this paper is neither to prove that the physical Universe globally has
the topology of $S^3$, nor to replace the standard cosmological model. The scope
is deliberately narrower: we construct a conditional, but mathematically
well-defined, reconstruction.

We take an observationally admissible small curvature modulus, together with the
closed-sign interpretation $\Omega_K^{S^3}<0$, not as a detected measurement but
as a fiducial benchmark for a conditional positive-curvature reconstruction. In
the standard FLRW convention,
\begin{equation}
  \Omega_K < 0
\end{equation}
corresponds to a closed, positively curved spatial geometry. If the spatial
section is modelled as a round three-sphere $S^3_R$, a large but finite curvature
radius is recovered from the small modulus $|\Omega_K|$.

This approach continues the discipline already introduced in our earlier
work~\cite{Preece2026causal}: $S^3$ is treated not as an arbitrary metaphor but
as the minimal compact simply connected spherical carrier in which the following
invariant structures meet: positive constant curvature, finite volume, a discrete
round spectrum, $SU(2)$ spin geometry, and the Hopf fibration. Parallelisability
alone is not sufficient to select $S^3$, since every closed orientable
three-manifold is parallelisable~\cite{BenedettiLisca2018}. The selection is
therefore made by the full invariant bundle, not by a single property. The
topology defines the carrier; dynamics, amplitudes, and physical fields require
additional laws.

In the companion paper on light and causal coherence~\cite{Preece2026causal},
the compact space $S^3_R \times \mathbb{R}$ is already used as the ambient model,
where the causal diameter equals
\begin{equation}
  D_c = \pi R,
\end{equation}
and the global causal accessibility time is
\begin{equation}
  t_{\rm coh} = \frac{\pi R}{c}.
\end{equation}
There it is emphasised that this is not an ``echo of light'', not superluminal
communication, and not a mystical interpretation, but finite null accessibility
within a compact, boundaryless space.

\section{Basic Reconstruction Hypothesis}
\label{sec:hypothesis}

Let the global spatial section be modelled as a round three-sphere:
\begin{equation}
  \Sigma_3 = S^3_R.
\end{equation}
Locally, for any sufficiently small observer domain:
\begin{equation}
  U \subset S^3_R, \qquad U \simeq \mathbb{R}^3.
\end{equation}
Therefore standard local physics is not abolished; it is preserved as a
quasi-local approximation.

Globally, however,
\begin{equation}
  S^3_R \neq \mathbb{R}^3.
\end{equation}
The differences are essential:
\begin{itemize}
  \item $S^3_R$ is compact, without boundary;
  \item it has finite proper volume;
  \item it has a finite causal diameter;
  \item its Laplacian has a discrete spectrum;
  \item its maximal geodesic distance is $\pi R$.
\end{itemize}

In this sense the flat $\mathbb{R}^3$ framework is not the unique ``default''
ontology. It is itself a choice: non-compactness, infinite volume, and the absence
of global closure. The compact $S^3$ framework offers a different choice: locally
nearly flat physics within a globally finite geometry.


\section{Invariant Geometry of the Round Carrier}
\label{sec:round-invariants}

For the round three-sphere of radius $R$,
\begin{equation}
  S^3_R=\{x\in\mathbb R^4:\norm{x}=R\},
\end{equation}
the intrinsic distance between two points is
\begin{equation}
  d(x,y)=R\arccos\!\left(\frac{x\cdot y}{R^2}\right).
\end{equation}
All local curvature invariants are fixed by the single radius $R$. The Riemann,
Ricci, and scalar curvatures are
\begin{align}
  R_{ijkl}
  &=\frac{1}{R^2}(g_{ik}g_{jl}-g_{il}g_{jk}), \\
  R_{ij}&=\frac{2}{R^2}g_{ij}, \\
  {}^{(3)}R&=\frac{6}{R^2}.
\end{align}
The round carrier is locally symmetric:
\begin{equation}
  \nabla_m R_{ijkl}=0.
\end{equation}
Consequently every local scalar curvature invariant is a polynomial in $R^{-2}$.
For example,
\begin{align}
  R_{ijkl}R^{ijkl}&=\frac{12}{R^4}, \\
  R_{ij}R^{ij}&=\frac{12}{R^4}, \\
  \left({}^{(3)}R\right)^2&=\frac{36}{R^4}.
\end{align}
These identities complete the local invariant description of the round carrier:
there is no hidden local curvature datum beyond the radius.

\subsection{No invariant spatial direction on round \texorpdfstring{$S^3$}{S3}}
\label{sec:no-invariant-direction}

The round simply connected carrier has no nonzero invariant spatial direction.
At any point $p\in S^3_R$, the stabilizer subgroup of the isometry group is
$SO(3)$, acting on the tangent space $T_pS^3$. This action is transitive on the
unit sphere in $T_pS^3$. If a tangent vector $v\in T_pS^3$ were invariant under
all rotations in the stabilizer, then
\begin{equation}
  h_\ast v=v\qquad \forall h\in SO(3),
\end{equation}
which is possible only for
\begin{equation}
  v=0.
\end{equation}
Thus a round $S^3_R$ geometry by itself cannot generate a preferred axis, a
preferred ring orientation, or a privileged spatial direction. In the same full
$SO(4)$ sense, no nonzero invariant one-form or two-form is selected by the bare
round carrier. Such features may arise only from data, observer position, survey
masks, field dynamics, initial conditions, a Hopf-frame choice, or a nontrivial
quotient topology $S^3/\Gamma$. The invariant content of the round carrier
consists of distances, volumes, curvature, spectra, and correlation functions,
not directions.

\section{Invariant Selection Principle for \texorpdfstring{$S^3$}{S3}}
\label{sec:invariant-selection}

The word ``unique'' is used in this paper only under a precise bundle of
assumptions. Positive curvature by itself does not determine global topology:
closed spherical space forms may be quotients $S^3/\Gamma$ by finite freely acting
isometry groups. Likewise, parallelisability by itself does not select $S^3$, since
all closed orientable three-manifolds are parallelisable~\cite{BenedettiLisca2018}.
The strict selection statement is therefore conditional:
\begin{quote}
\textbf{Invariant Selection Principle.}
If the spatial carrier is required to be a smooth, compact, boundaryless,
simply connected three-manifold with constant positive sectional curvature and
with the minimal spinorial Lie-group carrier structure required by $SU(2)$, then
the carrier is the round three-sphere $S^3_R$.
\end{quote}

The invariant support chain is:
\begin{align}
\text{closed + simply connected + dimension }3
  &\Rightarrow S^3
  \quad (\text{Poincar\'e--Perelman}), \\
K>0\text{ and simply connected}
  &\Rightarrow (S^3_R,g_R)\text{ as the spherical model}, \\
S^3 \simeq SU(2) \simeq Sp(1)
  &\Rightarrow \text{natural spin/rotation double-cover carrier}, \\
S^1 \hookrightarrow S^3 \to S^2,
  \qquad [S^3\to S^2]
  &\in\pi_3(S^2)\simeq\mathbb Z .
\end{align}
where the last line is the nontrivial Hopf fibration. The theorem proved by
Perelman, via Ricci flow and geometrization, states that every closed simply
connected three-manifold is a three-sphere~\cite{MorganTian2007,ClayPoincare}.
The identification $S^3\simeq SU(2)$ follows from unit quaternions and gives the
standard double-cover relation $SU(2)\to SO(3)$ relevant for spin
representations~\cite{SU2LieNotes}. The Hopf fibration supplies a nontrivial
$S^1$--fibered total space whose total space is $S^3$.

Thus $S^3$ is not claimed to be forced by observation alone. It is the unique
object selected when the above invariant requirements are imposed simultaneously.
This is the strict version of the topological claim used below.

\section{Comparison of Thurston Geometries}
\label{sec:thurston-comparison}

To illustrate the selection logic, Table~\ref{tab:thurston-compare} compares the
eight Thurston model geometries appearing in geometrization. The table does not
claim that no compact quotients exist in other geometries. Compact Euclidean and
hyperbolic quotients, for example, can have discrete Laplacian spectra. The point
is narrower: among the simply connected Thurston model carriers, only the
spherical model simultaneously supplies compact closure, simple connectivity,
constant positive sectional curvature, the $SU(2)$/Hopf structure, round spectral
data, and native integer topological charges without imposing boundary conditions
at infinity.
\begin{table}[htbp]
\centering
\caption{\label{tab:thurston-compare}Thurston model geometries under the invariant bundle used in this reconstruction. Compact quotients of other model geometries may have discrete spectra, but they are not simply connected model carriers.}
\scriptsize
\renewcommand{\arraystretch}{1.12}
\resizebox{0.98\textwidth}{!}{%
\begin{tabular}{llllll}
\toprule
Geometry & Model & $K$ & Compact SC? & $SU(2)$/Hopf? & Charges \\
\midrule
Spherical & $S^3$ & $+$ & Yes & Yes & Yes \\
Euclidean & $\mathbb E^3$ & $0$ & No & No & quotient \\
Hyperbolic & $\mathbb H^3$ & $-$ & No & No & quotient \\
$S^2\!\times\!\mathbb R$ & $S^2\!\times\!\mathbb R$ & mixed & No & partial & no \\
$\mathbb H^2\!\times\!\mathbb R$ & $\mathbb H^2\!\times\!\mathbb R$ & mixed & No & No & no \\
$\widetilde{SL(2,\mathbb R)}$ & universal cover & mixed & No & partial & no \\
Nil & Heisenberg & mixed & No & No & no \\
Sol & solvable & mixed & No & No & no \\
\bottomrule
\end{tabular}%
}
\end{table}

\subsection{Relation to the Poincar\'e Dodecahedral Space}
\label{sec:pds-comparison}

A useful comparison object is the Poincar\'e dodecahedral space,
\begin{equation}
  M_{\rm PDS}=S^3/I^\ast ,
\end{equation}
where $I^\ast$ is the binary icosahedral group of order $120$. This is a compact
multiply connected spherical space form of constant positive sectional curvature.
Its fundamental domain may be represented by a regular dodecahedron whose opposite
faces are identified after a $36^\circ$ twist. Since the quotient has $120$
copies in the simply connected cover, its volume at the same curvature radius is
\begin{equation}
  V_{\rm PDS}
  =
  \frac{1}{120}V_{S^3}
  =
  \frac{2\pi^2R^3}{120}.
\end{equation}
It is a homology three-sphere, but not a simply connected one:
\begin{equation}
  H_\ast(M_{\rm PDS})\simeq H_\ast(S^3),
  \qquad
  \pi_1(M_{\rm PDS})\simeq I^\ast\neq 0.
\end{equation}

The Poincar\'e dodecahedral space is historically important in CMB topology
because it was proposed as a finite positively curved model capable of suppressing
large-angle temperature correlations and reducing low multipole power~\cite{Luminet2003,Aurich2005}.
In such a quotient, the scalar harmonics are not the full set of round-$S^3$
harmonics but the subset invariant under the deck group $I^\ast$. Equivalently,
let $\mathcal H^{(n)}_{S^3}$ denote the scalar eigenspace of $-\Delta_{S^3}$
with eigenvalue $n(n+2)/R^2$. For each such eigenspace one keeps only
\begin{equation}
  \mathcal H^{(n)}_{\rm PDS}
  =
  \left\{
  Y\in\mathcal H^{(n)}_{S^3}
  :\rho(g)Y=Y\;\;\forall g\in I^\ast
  \right\},
\end{equation}
so the multiplicity is the dimension of the trivial $I^\ast$ representation
inside the restricted $S^3$ eigenspace. In character language this is the group
projection
\begin{equation}
  m^{\rm PDS}_n
  =
  \frac{1}{|I^\ast|}
  \sum_{g\in I^\ast}\chi_n(g),
\end{equation}
where $\chi_n$ is the character of the corresponding round-$S^3$ harmonic
representation restricted to $I^\ast$. Thus the PDS spectrum is a group-projected
subcatalogue of the $S^3$ spectrum, not a new local curvature spectrum.
Published PDS CMB calculations often use the notation
$\lambda_\beta=(\beta^2-1)/R^2$, related to the present convention by
\begin{equation}
  \beta=n+1,
  \qquad
  \lambda_n=\frac{n(n+2)}{R^2}=\frac{\beta^2-1}{R^2}.
\end{equation}
In this notation the first nonconstant surviving PDS scalar modes occur only at
comparatively high $\beta$ values, such as $\beta=13$~\cite{Aurich2005}. This
illustrates the stronger infrared gap and reduced low-mode degeneracy relative to
the simply connected cover. The radius $R$ and local curvature $K=1/R^2$ are the
same; the difference is the quotient projection on the allowed global modes.

The observational signature is also different. A multiply connected quotient may
produce repeated images of the last-scattering surface and hence pairs of matched
circles in the CMB sky~\cite{Cornish1998,PlanckTopology2016}. For the Poincar\'e
dodecahedral model the standard signature consists of six pairs of matched
circles with the characteristic dodecahedral phase shift. By contrast, the simply
connected sector adopted here,
\begin{equation}
  \Gamma=\{1\},
  \qquad
  M=S^3_R,
\end{equation}
has no nontrivial deck transformations and therefore no matched-circle signature
of this quotient type. The absence of detected matched circles in Planck rules out
some small fundamental domains but does not close the topology programme; current
COMPACT analyses emphasize that topology searches remain computationally and
statistically open~\cite{COMPACTPromise2024,COMPACTPartIIA2024}. Recent lens-space
circle searches also show how quotient sectors $L(p,q)=S^3/\mathbb Z_p$ are
constrained by the absence of matched circles, while the simply connected sector
has no such quotient-circle signature~\cite{COMPACTLens2025}.

This distinction clarifies the role of the Invariant Selection Principle. The
present paper does not claim that all spherical space forms are impossible.
Rather, it deliberately selects the minimal simply connected carrier. Curvature
alone fixes only the local spherical model and radius; it does not choose between
$S^3$ and quotients $S^3/\Gamma$. The additional simply connected requirement
sets $\Gamma=\{1\}$ and therefore selects the universal cover $S^3_R$.

Thus the Poincar\'e dodecahedral space is best understood here as the leading
excluded quotient sector. If future CMB temperature and polarization searches were
to detect the specific matched-circle or off-diagonal correlation pattern of a
nontrivial quotient, the simply connected sector choice would have to be revised.
If such signatures remain absent while low-$\ell$ anomalies persist, the
large-radius simply connected compact-mode programme remains the cleaner minimal
option.

For a general quotient $M=S^3/\Gamma$, the scalar eigenspace is
\begin{equation}
  \mathcal H_\beta(M)=
  \{Q\in\mathcal H_\beta(S^3):Q(gx)=Q(x)\;\forall g\in\Gamma\},
\end{equation}
with multiplicity
\begin{equation}
  m_\beta^\Gamma
  =\frac{1}{|\Gamma|}\sum_{g\in\Gamma}\chi_\beta(g).
\end{equation}
Such quotient sectors can generate off-diagonal CMB correlations of the form
\begin{equation}
  C_{\ell m,\ell'm'}^\Gamma
  =\sum_{\beta,s}\mathcal P_\beta
  \Delta_{\ell\beta}\Delta_{\ell'\beta}
  \xi_{\beta s;\ell m}\xi^*_{\beta s;\ell'm'} .
\end{equation}
For the simply connected round sector $\Gamma=\{1\}$, statistical isotropy gives
\begin{equation}
  C_{\ell m,\ell'm'}=C_\ell\delta_{\ell\ell'}\delta_{mm'}.
\end{equation}
Therefore any attempt to explain a persistent axis or orientation by topology
alone must move from the round simply connected carrier to a quotient sector or
to additional anisotropic dynamics.

\begin{table}[htbp]
\centering
\caption{Comparison between the simply connected round carrier used in this paper and the Poincar\'e dodecahedral quotient sector. At fixed curvature radius $R$, the quotient changes the allowed global modes and volume, not the local sectional curvature $K=1/R^2$.}
\label{tab:s3-pds-comparison}
\small
\renewcommand{\arraystretch}{1.16}
\resizebox{0.98\textwidth}{!}{%
\begin{tabular}{lll}
\toprule
Feature & Simply connected carrier $S^3_R$ & Poincar\'e dodecahedral space $S^3/I^\ast$ \\
\midrule
Topology & $\pi_1(S^3)=0$ & $\pi_1(S^3/I^\ast)\simeq I^\ast$, $|I^\ast|=120$ \\
Sector choice & Minimal universal cover, $\Gamma=\{1\}$ & Nontrivial spherical quotient, $\Gamma=I^\ast$ \\
Volume at same $R$ & $V_{S^3}=2\pi^2R^3$ & $V_{S^3}/120$ \\
Spectrum & Full round-$S^3$ hyperspherical spectrum & $I^\ast$-invariant subset of $S^3$ modes \\
Low-mode effect & Subtle compact-mode transfer effect for large $R$ & Stronger infrared gap from group projection at fixed $R$ \\
CMB signature & No quotient matched-circle signature & Six matched-circle pairs with dodecahedral phase relation \\
Status in this paper & Selected by the simply connected invariant bundle & Excluded quotient sector / falsification alternative \\
\bottomrule
\end{tabular}%
}
\end{table}

\section{Quantum-Scale Consequences of the Round Carrier}
\label{sec:quantum-s3-carrier}

Global compactness becomes directly measurable only on scales comparable to the
curvature radius $R_{\mathcal S^3}\approx548.32\,\mathrm{Gly}$. On all smaller
laboratory, atomic, nuclear, and ordinary cosmological scales inside the particle
horizon, every sufficiently small patch $U\subset S^3_R$ remains locally
Euclidean, $U\simeq\mathbb R^3$, so the local form of standard physics is
preserved. The quantum relevance of the round carrier is therefore not a local
``new force''. It appears through global spectral, spinorial, and topological
data that are intrinsic to round $S^3$ rather than to an arbitrary
three-manifold.

\subsection{Conservation: Strict Form and Limitation}

For a compact spatial section without boundary there is no surface at spatial
infinity through which a flux term can escape. In a static product spacetime
$S^3_R\times\mathbb R$ with a timelike Killing field $X=\partial_t$, the Noether
charge
\begin{equation}
  E_X[\Sigma]
  =\int_{\Sigma}\! T_{\mu\nu}X^\mu n^\nu\,d\Sigma
\end{equation}
is conserved under the usual stress-energy conservation law
$\nabla_\mu T^{\mu\nu}=0$. This is the correct strict statement. In a genuinely
dynamical FLRW spacetime there is generally no global timelike Killing vector,
so one should not claim an absolute globally conserved energy in the same sense;
the invariant statement is local covariant conservation plus exact global
charges only for actual symmetries.

Thus the role of $S^3$ is precise: compactness removes boundary-at-infinity
losses, while exact global conservation requires the corresponding Killing
symmetry. This correction prevents the argument from confusing compact topology
with time-translation invariance.


\subsection{Round Spectral Data}
\label{sec:round-spectral-data}

On the round three-sphere the scalar Laplace--Beltrami spectrum is discrete. It is
convenient to use the standard integer $\beta=n+1$:
\begin{equation}
  -\Delta_{S^3}Q_{\beta\ell m}
  =\frac{\beta^2-1}{R^2}Q_{\beta\ell m},
  \qquad
  \beta=1,2,3,\ldots .
\end{equation}
For each $\beta$,
\begin{equation}
  \ell=0,1,\ldots,\beta-1,
  \qquad
  m=-\ell,\ldots,\ell,
\end{equation}
and the scalar degeneracy is
\begin{equation}
  g_\beta=\beta^2.
\end{equation}
This is the same spectrum as
\begin{equation}
  -\Delta Y_{n\ell m}=\frac{n(n+2)}{R^2}Y_{n\ell m},
  \qquad g_n=(n+1)^2,
\end{equation}
with $\beta=n+1$. To avoid mixing conventions, the correspondence is
\begin{center}
\small
\renewcommand{\arraystretch}{1.12}
\begin{tabular}{llll}
\toprule
Old index & Hyperspherical index & Scalar eigenvalue & Degeneracy \\
\midrule
$n=0$ & $\beta=1$ & $0$ & $1$ \\
$n=1$ & $\beta=2$ & $3/R^2$ & $4$ \\
$n=2$ & $\beta=3$ & $8/R^2$ & $9$ \\
$n$ & $\beta=n+1$ & $n(n+2)/R^2=(\beta^2-1)/R^2$ & $(n+1)^2=\beta^2$ \\
\bottomrule
\end{tabular}
\end{center}

The addition theorem gives the spectral kernel in closed form:
\begin{align}
  \sum_{\ell m}Q_{\beta\ell m}(x)Q^*_{\beta\ell m}(y)
  &=\frac{\beta}{2\pi^2R^3}\frac{\sin(\beta\gamma)}{\sin\gamma}, \\
  \gamma&=\frac{d(x,y)}{R}.
\end{align}
This is the invariant bridge between compactness and two-point correlations.
A derivation of this addition theorem, its normalization, and its flat-space
limit is given in Appendix~\ref{app:addition-theorem}.
The heat trace is correspondingly a sum, not a continuum integral:
\begin{equation}
  Z(t)=\sum_{\beta=1}^{\infty}\beta^2
  \exp\!\left[-t\frac{\beta^2-1}{R^2}\right].
\end{equation}
The difference from $\mathbb R^3$ is therefore not a local force; it is the
replacement of a continuous $k$-integral by a discrete hyperspherical catalogue.

The spectrum of free fields on $S^3_R\times\mathbb R$ is a discrete mode problem
rather than an infinite-volume continuum problem. This gives exact
hyperspherical-harmonic expansions for scalar, vector, tensor, and spinorial
fields and is the mathematical basis of a compact-mode CMB and vacuum-energy
programme. Vacuum-energy sums are still infinite discrete sums and require a
renormalization prescription such as zeta-function, heat-kernel, or physical
subtraction. The improvement relative to $\mathbb R^3$ is the absence of an
infinite-volume infrared continuum and the presence of exact degeneracies and mode
gaps. Casimir calculations in Einstein-universe backgrounds illustrate precisely
this role of curvature and topology in renormalized vacuum energy~\cite{Mota2022}.

\subsection{Spinors and the Dirac Gap}

For the round spin three-sphere, the Dirac spectrum is explicitly known; with a
common sign convention one may write
\begin{equation}
  \lambda_k^{\pm}=\pm\frac{k+3/2}{R},
  \qquad
  k=0,1,2,\ldots .
\end{equation}
Thus the round geometry supplies a nonzero spectral gap of order $R^{-1}$ for
the untwisted Dirac operator~\cite{CamporesiHiguchi1996}. The Lichnerowicz
formula implies that a compact spin manifold with positive scalar curvature has
no untwisted harmonic spinors~\cite{GromovLawson1983}. This is a strong stability
statement for the simplest spin sector.

The limitation is equally important: twisted Dirac operators, gauge backgrounds,
or additional bundles can alter zero-mode questions. Therefore the article uses
the Dirac gap as a controlled round-$S^3$ spectral fact, not as a universal proof
that all fermionic sectors are automatically anomaly-free.

\subsection{Hopf Fibration and Quantum Phase Geometry}

A normalized one-qubit state is a point on $S^3\subset\mathbb C^2$ before
quotienting by the physically irrelevant global phase. The Hopf fibration
\begin{equation}
  S^1\hookrightarrow S^3\twoheadrightarrow S^2
\end{equation}
then sends the total state sphere to the Bloch sphere, while the circular fiber
records the overall phase degree of freedom. This is the standard geometric
relation between qubit state geometry, projective Hilbert space, and phase
holonomy~\cite{MosseriDandoloff2001}.

For field configurations the same domain $S^3$ supports the Hopf invariant:
\begin{equation}
  [S^3,S^2]\simeq\pi_3(S^2)\simeq\mathbb Z .
\end{equation}
The linking of inverse images under maps $S^3\to S^2$ is the geometric content of
this invariant, and it underlies Hopfions and knotted field configurations in
Faddeev--Skyrme type models~\cite{Urbantke2003,Sutcliffe2007}. On ordinary
$\mathbb R^3$ this structure is normally obtained only after imposing a boundary
condition at infinity and compactifying $\mathbb R^3\cup\{\infty\}\simeq S^3$.
In a genuine compact $S^3$ carrier the same topology is native rather than
artificially added.

\subsection{Skyrmion Degree and Native Baryon Number}

For the Skyrme field $U$, the target is $SU(2)\simeq S^3$. The topological
charge is the degree of a map into $SU(2)$; in the usual $\mathbb R^3$ treatment
this is well-defined after a boundary condition at infinity, effectively turning
space into $S^3$~\cite{Manton1998}. On a compact $S^3$ spatial carrier the same
topological degree is intrinsic:
\begin{equation}
  U:S^3\to SU(2)\simeq S^3,
  \qquad
  B=\deg(U)\in\mathbb Z .
\end{equation}
This supplies a precise sense in which a compact $S^3$ carrier makes Skyrmion
number global without imposing an external point at infinity.

\begin{center}
\small
\renewcommand{\arraystretch}{1.18}
\resizebox{0.98\textwidth}{!}{%
\begin{tabular}{lll}
\toprule
Claim type & Strict invariant content & Required wording in this article \\
\midrule
Local physics & $U\subset S^3_R$, $U\simeq\mathbb R^3$ & Standard local physics is preserved. \\
Global energy & No boundary-at-infinity flux; exact energy needs $\partial_t$ Killing symmetry & Static conservation; local covariant conservation in FLRW. \\
Vacuum energy & Discrete spectrum and finite volume; UV sums still need renormalization & Computable renormalized mode sums, not finite bare sums. \\
Dirac sector & Positive scalar curvature; explicit round-$S^3$ gap & Untwisted spectral stability; twisted sectors separate. \\
Hopf/qubit sector & $S^1\hookrightarrow S^3\to S^2$, $\pi_3(S^2)=\mathbb Z$ & Native phase, bundle, and Hopfion topology. \\
Skyrmion sector & $S^3\to SU(2)\simeq S^3$, degree in $\mathbb Z$ & Native global baryon-number topology. \\
Uniqueness & Full invariant bundle, not compactness alone & Conditional selection, not observational proof. \\
\bottomrule
\end{tabular}%
}
\end{center}

This section strengthens the article without overclaiming. The round $S^3$
carrier is not required in order to recover local quantum mechanics; it is
selected by the demand that the global carrier simultaneously provide compact
closure, exact round spectra, Hopf topology, $SU(2)$ spin geometry, and native
integer topological charges.

\section{Reconstructive Curvature Principle}
\label{sec:rcp}

The logical bridge from observation to geometry is formalised as follows.

\begin{quote}
\textbf{Reconstructive Curvature Principle (RCP).}
\textit{If the spatial curvature parameter satisfies $\Omega_K < 0$ within
observational bounds, and if this value is treated not as noise but as a
conditional signal of positive spatial curvature, then the standard FLRW
curvature radius formula
\begin{equation}
  R_c = \frac{c/H_0}{\sqrt{|\Omega_K|}}
  \label{eq:rcp}
\end{equation}
specifies the simply connected constant-positive-curvature universal cover
$S^3_{R_c}$. The physical global section may in principle be a spherical space
form $S^3/\Gamma$; the minimal simply connected sector sets $\Gamma=\{1\}$ and
therefore selects $S^3_{R_c}$ itself.}
\end{quote}

Four points of discipline govern this principle:

\begin{enumerate}
\item \textbf{The RCP is conditional, not demonstrative.}
  It does not prove $\Omega_K < 0$ is real; it constructs what follows
  \emph{if} it is.

\item \textbf{The sign is decisive.}
  The FLRW convention assigns $\Omega_K < 0$ to positive curvature (closed),
  $\Omega_K > 0$ to negative curvature (open), and $\Omega_K = 0$ to flat.
  A positive curvature reconstruction is therefore active only when
  $\Omega_K < 0$ is robustly non-zero.

\item \textbf{The formula is exact within the model.}
  Given $H_0$ and $\Omega_K$, equation~\eqref{eq:rcp} is algebraically exact.
  Uncertainty enters only through the observational inputs.

\item \textbf{Topology requires a separate sector choice.}
  Curvature fixes the local spherical model and radius. It does not by itself
  exclude spherical quotient topologies $S^3/\Gamma$. The present paper uses the
  minimal simply connected carrier sector.
\end{enumerate}

For the canonical Light-Trisphere reconstruction adopted in this paper, the
input is not the Planck-only closed-universe anomaly. It is the fiducial
radius-modulus benchmark
\begin{equation}
\begin{aligned}
  |\Omega_K| &= 0.0007, \\
  \Omega_K^{S^3} &= -0.0007, \\
  H_0 &= 67.4\,\mathrm{km\,s}^{-1}\mathrm{Mpc}^{-1}.
\end{aligned}
\label{eq:canonical_inputs}
\end{equation}
The sign convention is essential: in standard FLRW notation $\Omega_K<0$ means
closed positive spatial curvature. Previous curvature-modulus estimates quoting
$\Omega_K=+0.0007$ are therefore used here only through the same modulus
$|\Omega_K|=0.0007$, with the sign flipped for the closed $S^3$ interpretation.
With these inputs,
\begin{equation}
  R_c = \frac{c/H_0}{\sqrt{|\Omega_K|}}
      = 548.3\Gly,
\end{equation}
yielding the fiducial radius adopted throughout this paper:
\begin{equation}
  \boxed{R_{\mathcal{S}^3} = 548.3\Gly.}
\end{equation}

\subsection{Degeneration at Zero Curvature}

The limiting behaviour as $|\Omega_K| \to 0$ shows why a small non-zero value
is categorically different from zero:
\begin{equation}
\begin{aligned}
  |\Omega_K| \to 0
  &\Rightarrow R_c \to \infty \\
  &\Rightarrow S^3_{R_c} \to \mathbb{R}^3
  \quad\text{(locally and in the limit).}
\end{aligned}
\end{equation}
The closed spherical model decompactifies into the flat limit. Conversely, even a
tiny non-zero closed-sign curvature modulus produces a \emph{finite} spherical
curvature radius:
\begin{equation}
\begin{aligned}
  \Omega_K < 0
  &\Rightarrow R_c < \infty \\
  &\Rightarrow S^3_{R_c}\text{ is compact}
  \quad\text{(simply connected sector).}
\end{aligned}
\end{equation}
For $\Omega_K>0$ the corresponding FLRW sector is open/negative-curvature and
belongs to the hyperbolic side, not to the $S^3$ reconstruction. The full
reconstruction chain from observation to observable consequence is:
\begin{equation}
\begin{aligned}
  \Omega_K &\mapsto R_c \mapsto S^3_{R_c} \\
           &\mapsto D_{\rm obs} \\
           &\mapsto \text{modes, causal relations, horizons, structures.}
\end{aligned}
\end{equation}
A small closed-sign curvature is not ``nearly zero'' in the geometric sense; it
fixes a finite radius for the spherical model. Zero gives decompactification; a
closed-sign non-zero value gives a distinct compact spherical carrier in the
chosen simply connected sector, with a specific spectrum, volume, and causal
structure.


\section{Curvature-Distance Law and Reconstruction Test}
\label{sec:distance-law}

The fiducial radius is only the kinematic starting point. Observable predictions
require the FLRW distance-redshift relation, perturbation transfer functions, and
a likelihood comparison to data. Let
\begin{equation}
  E(z)=\frac{H(z)}{H_0}.
\end{equation}
The line-of-sight comoving distance is
\begin{equation}
  D_C(z)=\frac{c}{H_0}\int_0^z\frac{dz'}{E(z')}.
\end{equation}
For a closed spatial section with $\Omega_K<0$, the dimensionless radial coordinate
is
\begin{equation}
  \chi(z)=\frac{D_C(z)}{R_c}
  =\sqrt{|\Omega_K|}\int_0^z\frac{dz'}{E(z')}.
\end{equation}
The transverse comoving distance is therefore
\begin{equation}
\boxed{
  D_M(z)=R_c\sin\chi(z)
  =\frac{c/H_0}{\sqrt{|\Omega_K|}}
  \sin\!\left[
    \sqrt{|\Omega_K|}\int_0^z\frac{dz'}{E(z')}
  \right].}
\end{equation}
The usual angular-diameter and luminosity distances are
\begin{equation}
  D_A(z)=\frac{D_M(z)}{1+z},
  \qquad
  D_L(z)=(1+z)D_M(z),
\end{equation}
and the differential comoving volume per redshift and solid angle is
\begin{equation}
  \frac{dV}{dz\,d\Omega}=\frac{c\,D_M^2(z)}{H(z)}.
\end{equation}
For small curvature the expansion reads
\begin{equation}
  D_M=D_C\left[
  1+\frac{\Omega_K}{6}\left(\frac{H_0D_C}{c}\right)^2
  +O(\Omega_K^2)\right].
\end{equation}
Thus a closed model with $\Omega_K<0$ gives a negative leading correction to the
flat transverse distance at fixed $D_C$. At the particle-horizon scale used in this
paper, $D_C\simeq46.13\,\mathrm{Gly}$, the fiducial value
$\Omega_K=-0.0007$ gives the order-of-magnitude diagnostic
\begin{equation}
  \frac{D_M-D_C}{D_C}
  \simeq -1.179\times10^{-3},
\end{equation}
which is observable only through a full joint distance likelihood, not from the
geometry alone.

A model-independent FLRW consistency reconstruction is
\begin{equation}
\boxed{
  \Omega_K(z)=
  \frac{[H(z)D_M'(z)]^2-c^2}{H_0^2D_M^2(z)} .}
\label{eq:ok_reconstruction}
\end{equation}
For any FLRW geometry this quantity must be constant in redshift. The closed
$S^3$ reconstruction requires, within observational errors,
\begin{equation}
  \Omega_K(z)<0,
  \qquad
  \frac{d\Omega_K}{dz}=0.
\end{equation}
This is a stronger observational test than selecting a single fiducial curvature
modulus. In current data the curvature benchmark used here remains a permitted
modulus rather than a detection: Planck 2018 combined with BAO gives a curvature
constraint consistent with flatness~\cite{Planck2020}, while DESI DR2 provides
high-precision BAO distances from more than $14$ million galaxies and quasars and
primarily sharpens the distance-redshift likelihood rather than proving positive
spatial curvature~\cite{DESIDR2}. ACT DR6 extended-model constraints similarly
report no statistically significant preference for departures from baseline
$\Lambda$CDM, reinforcing the need to treat the present $S^3_R$ construction as a
registered likelihood programme rather than as a detected curvature sector~\cite{ACTDR6Extensions}.

\subsection{BAO and acoustic-scale diagnostic observables}
\label{sec:bao-acoustic-diagnostics}

The distance law enters present and future BAO analyses through the invariant
combinations
\begin{align}
  D_H(z)&=\frac{c}{H(z)},\qquad
  D_V(z)=\left[zD_H(z)D_M^2(z)\right]^{1/3},\\
  F_{\rm AP}(z)&=\frac{D_M(z)H(z)}{c},\qquad
  \ell_A=\pi\frac{D_M(z_\ast)}{r_s(z_\ast)}.
\end{align}
At fixed line-of-sight distance $D_C$, the leading closed-curvature correction is
\begin{equation}
  \frac{D_M-D_C}{D_C}
  =\frac{\Omega_K}{6}\left(\frac{H_0D_C}{c}\right)^2
  +O(\Omega_K^2),
  \qquad \Omega_K<0.
\end{equation}
Thus the purely geometric transverse correction is negative for the closed sector.
For volume and acoustic observables, however, the measured shift also depends on
which background expansion history $E(z)$ is compared. Therefore
this paper does not register a standalone BAO number as a detection. The
registered observable is the joint likelihood of
\begin{equation}
  \left\{D_M/r_d,\ D_H/r_d,\ D_V/r_d,\ F_{\rm AP},\ \ell_A\right\}
\end{equation}
under the same priors and covariance as the flat baseline.

\subsection{Compact mode-shift diagnostic}
\label{sec:mode-shift-diagnostic}

The compact spectrum samples wavenumbers at
\begin{equation}
  k_\beta=\frac{\sqrt{\beta^2-1}}{R},
\end{equation}
rather than integrating continuously over $k$. For a smooth transfer integrand,
the leading effect of replacing an integral by a sum may be diagnosed by the local
spectral slope rather than asserted as a universal quadrupole suppression:
\begin{equation}
\begin{aligned}
  \Delta C_\ell^{XY}\sim
  &\left.\frac{d C_\ell^{XY}}{d\ln k}\right|_{k=k_\ell}
  \Delta\ln k_\ell \\
  &+\text{window/covariance terms}.
\end{aligned}
\end{equation}
This expression is only a diagnostic expansion. The prediction used for data
analysis remains the full hyperspherical transfer sum of
Sec.~\ref{sec:cmb-low-l} and Sec.~\ref{sec:transfer-likelihood-compressed}.


\section{Scope Guardrail: Background Expansion and Effective Acceleration}
\label{sec:background-expansion-scope}

The reconstruction developed in this paper is geometric and spectral. It fixes a
compact carrier, a curvature radius, the associated distance law, and a compact
mode registry. It does not determine the physical origin of late-time acceleration.
Consequently no statement in this article should be read as a proof that dark
energy or dark matter do not exist, or as a derivation of a replacement for the
late-time background sector.

Any effective acceleration component must enter only through a specified expansion
history,
\begin{equation}
  E(z)=\frac{H(z)}{H_0},
\end{equation}
and must be tested through the same BAO, supernova, CMB, lensing, and growth
likelihoods used for standard cosmological models. The present work therefore
uses $E(z)$ as an external background input for the distance and transfer-function
pipeline. It does not import a separate dark-sector theory into the main
reconstruction.

\section{Numerical Anchoring}
\label{sec:numerical}

In the main text we quote observationally meaningful rounded values. Machine-precision
numbers are retained only as fiducial benchmark outputs in Table~\ref{tab:summary}.
For the adopted modulus one obtains
\begin{equation}
  R_{\mathcal{S}^3}\simeq 548.3\Gly,
\end{equation}
where $R_{\mathcal{S}^3}$ is the curvature radius of the global three-sphere.

In units of $\mathrm{Gly}/\mathrm{Gyr}$, the speed of light is numerically unity:
\begin{equation}
  c = 1\,\mathrm{Gly/Gyr}.
\end{equation}
Hence the light-travel time corresponding to the curvature radius is:
\begin{equation}
  t_{\mathcal{S}^3} = \frac{R_{\mathcal{S}^3}}{c}
  \simeq 548.3\Gyr.
\end{equation}

\textbf{Important:}
\begin{equation}
  \boxed{t_{\mathcal{S}^3} \text{ is \emph{not} the age of the Universe.}}
\end{equation}
This is the light-travel time of the curvature radius. The age of the manifested
hot cosmological history remains a separate quantity:
\begin{equation}
  t_0 \simeq 13.79\Gyr.
\end{equation}

\section{Full Geometry of the Three-Sphere}
\label{sec:geometry}

\subsection{Diameter in the Embedding Space}

If $S^3_R$ is regarded as the standard sphere in $\mathbb{R}^4$,
\begin{equation}
  S^3_R = \{x \in \mathbb{R}^4 : |x| = R\},
\end{equation}
the diameter in the ambient $\mathbb{R}^4$ embedding is:
\begin{equation}
  \boxed{D_{\rm embed} = 2R = 1096.6\Gly.}
\end{equation}
This is \emph{not} the maximal intrinsic distance along the three-sphere itself.

\subsection{Intrinsic Causal Diameter \texorpdfstring{$D_{\rm causal} = \pi R$}{Dcausal = pi R}}

Two distinct diameter concepts must be carefully distinguished on $S^3_R$:
\begin{align}
  D_{\rm embed}   &= 2R    \quad \text{(diameter in ambient }\mathbb{R}^4\text{),} \\
  D_{\rm causal}  &= \pi R \quad \text{(maximal intrinsic geodesic distance).}
\end{align}
These differ by the factor $\pi/2 \approx 1.571$; conflating them leads to a
systematic overestimate of the intrinsic scale by $57\%$.

On $S^3_R$, the maximum intrinsic geodesic distance between antipodal points is:
\begin{equation}
  D_{\rm causal} \equiv D_{\rm geo,max} = \pi R.
\end{equation}
Numerically:
\begin{equation}
  \boxed{D_{\rm causal} = 1722.6\Gly.}
\end{equation}
This is the causal diameter: the maximum finite distance any null signal can
traverse within the compact spatial section. It is not a horizon in the
particle-horizon sense; it is the intrinsic diameter of the manifold.

\subsection{Full Great Circle}

The complete traversal along a great geodesic circle is
\begin{equation}
  C = 2\pi R.
\end{equation}
Numerically:
\begin{equation}
  \boxed{C_{\mathcal{S}^3} = 3445.2\Gly.}
\end{equation}
The corresponding light-travel time for a full traversal is
\begin{equation}
  t_{\rm loop} = \frac{2\pi R}{c} = 3445.2\Gyr.
\end{equation}
Half and quarter great circles are
\begin{equation}
\begin{aligned}
  \pi R &= 1722.6\Gly, \\
  \frac{\pi R}{2} &= 861.3\Gly .
\end{aligned}
\end{equation}

\section{Proper Volume of \texorpdfstring{$S^3$}{S3}}
\label{sec:volume}

The proper three-dimensional volume of the round three-sphere is:
\begin{equation}
  V_{S^3} = 2\pi^2 R^3.
\end{equation}
Substituting $R = 548.3\Gly$ gives
\begin{equation}
  \boxed{V_{\mathcal{S}^3} = 3.254\times 10^{9}\Gly^3.}
\end{equation}
This is \emph{not} the volume of an ordinary ball,
$V_{\rm ball} = 4\pi R^3/3$. It is the proper volume of the closed
three-dimensional manifold $S^3$ itself.

\section{Curvature}
\label{sec:curvature}

For the round three-sphere, the sectional curvature is:
\begin{equation}
  K_{\mathcal{S}^3} = \frac{1}{R^2}.
\end{equation}
Numerically:
\begin{equation}
  \boxed{K_{\mathcal{S}^3} = 3.326\times 10^{-6}\Gly^{-2}.}
\end{equation}
The three-dimensional scalar curvature is
\begin{equation}
  {}^{(3)}R = \frac{6}{R^2}.
\end{equation}
Numerically:
\begin{equation}
  \boxed{{}^{(3)}R = 1.996\times 10^{-5}\Gly^{-2}.}
\end{equation}
In SI units:
\begin{align}
  \boxed{|K|} &= 3.716\times 10^{-56}\,\mathrm{m}^{-2}, \\
  \boxed{|{}^{(3)}R|} &= 2.230\times 10^{-55}\,\mathrm{m}^{-2}.
\end{align}

The curvature is positive but extremely small. A quasi-local observer will
therefore consistently reconstruct their domain as nearly flat.

\section{Static-Carrier Age Diagnostic}
\label{sec:arc}

The age of the hot cosmological history is
\begin{equation}
  t_0 \simeq 13.79\Gyr.
\end{equation}
Its ratio to the light-travel time of the curvature radius is
\begin{equation}
  \frac{t_0}{t_{\mathcal{S}^3}}
  = \frac{13.788}{548.3}.
\end{equation}
Evaluating:
\begin{equation}
  \boxed{\frac{t_0}{t_{\mathcal{S}^3}} = 0.02515
  \approx 2.515\%.}
\end{equation}
The inverse ratio is
\begin{equation}
  \boxed{\frac{t_{\mathcal{S}^3}}{t_0} = 39.77,}
\end{equation}
i.e. the curvature-radius light-time exceeds the age of the observable hot
history by a factor of approximately $39.77$.

Expressing $ct_0$ as an arc on radius $R_{\mathcal{S}^3}$:
\begin{equation}
  \theta_0 = \frac{ct_0}{R_{\mathcal{S}^3}}
  = \frac{t_0}{t_{\mathcal{S}^3}}
  = 0.02515\,\mathrm{rad}.
\end{equation}
In degrees:
\begin{equation}
  \boxed{\theta_0 = 1.441^\circ.}
\end{equation}
Thus the static comparison length $ct_0$ corresponds to approximately
$1.44^\circ$ on a great circle of the fiducial carrier.

\textbf{Note:} This is a static-carrier diagnostic only. It is not the radius of
the observable Universe and not the FLRW causal light-cone distance. In FLRW,
causal propagation is controlled by the conformal distance
$\int c\,d\tau/a(\tau)$, while the particle horizon is larger than $ct_0$ because
of expansion.

\section{Observable Horizon within the Three-Sphere}
\label{sec:horizon}

For the benchmark observable-domain calculation we adopt the particle-horizon radius
\begin{equation}
  D_{\rm particle} = 46.13\Gly.
\end{equation}
This value is not a theorem of $S^3$ geometry. It is a background-dependent input fixed
by the chosen expansion history $E(z)$ and by the assumed radiation, matter, neutrino,
and curvature content used to compute the conformal particle horizon.
Its fraction of the three-sphere radius is
\begin{equation}
  \frac{D_{\rm particle}}{R_{\mathcal{S}^3}} = 0.08412.
\end{equation}
The corresponding angular size is
\begin{equation}
  \theta_{\rm particle}
  = \frac{D_{\rm particle}}{R_{\mathcal{S}^3}} \cdot \frac{180^\circ}{\pi}.
\end{equation}
Numerically:
\begin{equation}
  \boxed{\theta_{\rm particle} = 4.820^\circ.}
\end{equation}
Hence, the particle-horizon radius subtends approximately $4.82^\circ$ on the
great circle of $S^3$.

The diameter of the observable region is
\begin{equation}
  2D_{\rm particle} = 92.25\Gly.
\end{equation}
Its fraction of the full great circle is
\begin{equation}
  \frac{2D_{\rm particle}}{2\pi R_{\mathcal{S}^3}} = 0.02678
  \approx 2.678\%.
\end{equation}

The central result of this section is that the observable Universe can be a
small quasi-local domain of a large compact $S^3$.

\subsection{Volumetric Fraction of the Observable Domain}

The angular subtended arc gives an intuitive picture, but the volumetric fraction
is more quantitatively revealing. Define the geodesic half-angle:
\begin{equation}
  \chi = \frac{D_{\rm particle}}{R_{\mathcal{S}^3}}
       = 0.08412\,\mathrm{rad}.
\end{equation}
The fraction of the total $S^3$ volume enclosed within geodesic distance
$D_{\rm particle}$ of the observer follows from the $S^3$ volume element:
\begin{equation}
  \Gamma(\chi)
  = \frac{\chi - \tfrac{1}{2}\sin(2\chi)}{\pi}.
  \label{eq:Gamma_chi}
\end{equation}
This follows directly from the round-$S^3$ volume element
$dV=R^3\sin^2\chi\,d\chi\,d\Omega_2$, since
\begin{equation}
  V(\chi)=4\pi R^3\int_0^\chi\sin^2 u\,du
  =2\pi R^3\left(\chi-\tfrac12\sin 2\chi\right),
\end{equation}
and $V_{S^3}=2\pi^2R^3$. For small $\chi \ll 1$ (the regime relevant here),
the leading-order expansion is:
\begin{equation}
  \Gamma(\chi) \approx \frac{2}{3\pi}\,\chi^3
  \qquad (\chi \ll 1).
  \label{eq:Gamma_approx}
\end{equation}
Equation~\eqref{eq:Gamma_approx} shows that the volumetric fraction falls as the
\emph{cube} of the angular scale, making it far smaller than the linear angular
comparison ($4.82^\circ$) suggests. Substituting $\chi \simeq 0.08412$:
\begin{equation}
\begin{aligned}
  \Gamma_{\rm obs}
  &= \Gamma(\chi_{\rm particle}) \\
  &= 1.261\,377\,743\,887\,7\times10^{-4} \\
  &= 0.012\,613\,777\,438\,877\%.
\end{aligned}
\end{equation}
The small-angle approximation gives $0.012\,631\,642\%$, differing from the exact
value by about $0.14\%$ relative. Thus the observable Universe occupies about
$\mathbf{1.26\times10^{-2}\%}$ of
the total volume of $S^3_R$---not only angularly small, but volumetrically
negligible. The corrected value follows directly from the cubic small-angle
scaling in Eq.~\eqref{eq:Gamma_approx}; the linear angular fraction must not be
read as a volume fraction.

The $\Gamma(\chi)$ function is the spatial analogue of the causal fraction
$\Gamma(T)$ introduced in Sec.~\ref{sec:causal}; together they express the
incompleteness of the observer's domain in both space and time. The joint benchmark
reported in Table~I is defined exactly as
\begin{equation}
  \Gamma_{\rm ST}
  =\Gamma\!\left(\frac{ct_0}{R}\right)\Gamma(\chi_{\rm particle})
  =4.2551\times10^{-10}.
\end{equation}
This is a static-carrier diagnostic product, not a new causal horizon in dynamical
FLRW spacetime:
\begin{center}
\fbox{\begin{minipage}{0.92\columnwidth}
\centering
We observe a small quasi-local sector, not the global $S^3$.
\end{minipage}}
\end{center}

\section{The Quasi-Local Observer Domain}
\label{sec:quasilocal}

\subsection{Formal Definition and Strict Incompleteness}

The quasi-local observer domain is not identified with the global three-sphere. It
is defined by finite causal accessibility, the particle horizon, the age of the
manifested history, and the reconstruction procedure. Therefore, even under the
compact $S^3$ ontology, the observable region remains a partial perspective.

Formally, the causal domain of an observer at point $p$ after cosmic time $T$ is:
\begin{equation}
  D_{\rm obs}(p,T)
  = \{q \in S^3_R : d_{S^3}(p,q) \leq cT\},
\end{equation}
so that $D_{\rm obs}(p,T) \subseteq S^3_R$. The strict incompleteness theorem
holds as follows:
\begin{equation}
\begin{aligned}
  D_{\rm obs}(p,T) &\subset S^3_R, \\
  D_{\rm obs}(p,T) &\neq S^3_R.
\end{aligned}
\end{equation}
This holds for every finite manifested reconstruction. This is not a contingent
observational limitation. It is a structural consequence: a domain is defined by
its observer and its horizon. Were it to coincide with the full $S^3_R$, it would
cease to be \emph{quasi}-local and would instead be the global spatial section
itself---a categorically different object.

The fractional causal accessibility volume is
\begin{equation}
  \Gamma(T)
  = \frac{\mathrm{Vol}(D_{\rm obs}(p,T))}{\mathrm{Vol}(S^3_R)}.
\end{equation}
For all $T < t_{\rm coh}$,
\begin{equation}
  \Gamma(T) < 1.
\end{equation}

\subsection{Boundary Without Boundary}
\label{sec:boundary-without-boundary}

The finite curvature radius recovered by the closed-branch reconstruction is not
an edge of the Universe. It is a radius of curvature of the carrier geometry:
\begin{equation}
  R_c=\frac{c/H_0}{\sqrt{|\Omega_K|}},
  \qquad
  \Sigma\simeq S^3_{R_c}
\end{equation}
for the simply connected closed sector.
The global spatial carrier is compact and finite, but boundaryless:
\begin{equation}
  \partial S^3_R=\varnothing .
\end{equation}
By contrast, an observer-defined domain may have a nonempty relational boundary,
for example
\begin{equation}
  D_{\rm obs}\subset S^3_R,
  \qquad
  \partial D_{\rm obs}\simeq S^2 .
\end{equation}
Thus
\begin{equation}
  \boxed{
  \partial S^3_R=\varnothing,
  \qquad
  \partial D_{\rm obs}\neq\varnothing,
  \qquad
  \partial D_{\rm obs}\not\equiv\partial S^3_R .}
\end{equation}
The boundary of $D_{\rm obs}$ is a boundary of reconstruction, causal access,
windowing, or imposed boundary data; it is not an ontological boundary of space.

This distinction is also the correct gravitational one. In covariant phase-space
language the on-shell diffeomorphism current is exact, schematically
\begin{equation}
  J_\xi\approx dQ_\xi,
\end{equation}
so Hamiltonian generators for diffeomorphisms reduce to boundary functionals
rather than bulk charges~\cite{IyerWald1994}. For a compact boundaryless global
slice there is no spatial edge term. For a proper subregion $D\subset S^3_R$,
however, the quasi-local generator is encoded on $\partial D$:
\begin{equation}
  H_\xi[D]
  =
  \int_{\partial D}
  (Q_\xi-\xi\cdot B) .
\end{equation}
This is compatible with Brown--York type quasi-local energy constructions, but the
present paper does not evaluate such charges or introduce a York-time dynamics
sector~\cite{BrownYork1993}. The point used here is only categorical:
\begin{center}
\fbox{\begin{minipage}{0.92\columnwidth}
\centering
The universe has no global edge; observer domains may have quasi-local boundaries.
\end{minipage}}
\end{center}

\subsection{Reachability Is Not Empirical Reconstruction}

At $T = t_{\rm coh} = \pi R/c$, the condition $\Gamma(T) = 1$ is achieved---but
only in the idealised static model sense of \emph{null accessibility}:
\begin{equation}
\begin{aligned}
  \Gamma(t_{\rm coh}) = 1
  \quad &\Leftrightarrow \quad
  \text{every point of }S^3_R \\
  &\text{is null-reachable from }p .
\end{aligned}
\end{equation}
This is a statement about geometric \emph{reachability}---the existence of a null
geodesic connecting any two points---not about empirical \emph{reconstruction}.
Reachability does not imply that an observer has access to, has measured, or has
integrated the information from the entire $S^3_R$.

The distinction is essential:
\begin{align}
  \Gamma(t_{\rm coh}) = 1
  &\quad \Leftrightarrow \quad
    \text{global null accessibility (geometric fact)}, \\
  D_{\rm obs}(p, t_{\rm coh}) = S^3_R
  &\quad \not\Leftrightarrow \quad
    \text{full empirical reconstruction (physical claim)}.
\end{align}
The local null-cone structure is preserved at all times, and no single observer
instantaneously integrates the entire compact section. The $\Gamma(T)$ function
measures causal accessibility, not the actual data held by the observer.

\subsection{Protection against False Universalism}

The two key methodological principles of this framework are:
\begin{quote}
The quasi-local domain is the observational shadow of $S^3$, not $S^3$ itself.

$S^3$ is not exhausted by any quasi-local perspective.
\end{quote}

Standard flat cosmology performs the substitution
$\text{locally reconstructed near-flatness} \to \mathbb{R}^3\text{-ontology}$.
This programme refuses that substitution:
\begin{equation}
\begin{aligned}
  \text{local near-flatness}
  &\not\Rightarrow \mathbb{R}^3, \\
  &\text{nor non-compact global geometry.}
\end{aligned}
\end{equation}
Local near-flatness is compatible with any large compact carrier:
$U \subset S^3_R$, $U \simeq \mathbb{R}^3$, provided $R$ is sufficiently large.
The inference from local flatness to global non-compactness is an \emph{additional
ontological assumption}, not a deduction.

\section{Apparent Flatness and Curvature Sagitta Visualization}
\label{sec:flatness}

This section is an embedding visualization and curvature-intuition diagnostic only. It
is not an invariant observational curvature estimator and does not replace the FLRW
distance test in Eq.~\eqref{eq:ok_reconstruction}. To estimate the deviation from a
flat chord, one may use the sagitta (arc-to-chord deviation):
\begin{equation}
  s(L)=R-\sqrt{R^2-\left(\frac{L}{2}\right)^2},
\end{equation}
which for small $L/R$ reduces to
\begin{equation}
  s(L) \approx \frac{L^2}{8R}.
\end{equation}

\textbf{For the light-travel age $L = ct_0 = 13.788\Gly$:}
\begin{equation}
  \boxed{s = 0.043\,338\,547\Gly = 43.34\Mly.}
\end{equation}

\textbf{For the particle horizon $L \simeq 46.13\Gly$:}
\begin{equation}
  \boxed{s = 0.485\,222\,026\Gly.}
\end{equation}

\textbf{For the diameter of the observable Universe $L = 92.250\Gly$:}
\begin{equation}
  \boxed{s = 1.943\,473\,550\Gly.}
\end{equation}

These figures demonstrate: on the scale of the entire observable region, the
curvature is geometrically non-zero, but it is extraordinarily gentle. Local
reconstructions will therefore robustly yield a nearly flat picture.

\section{Causal Accessibility}
\label{sec:causal}

In the compact $S^3_R$ geometry, the maximum intrinsic distance is $D_c = \pi R$,
and the time of full causal accessibility is
\begin{equation}
  t_{\rm coh} = \frac{\pi R}{c}.
\end{equation}
For the chosen radius:
\begin{equation}
  \boxed{t_{\rm coh} = 1722.6\Gyr.}
\end{equation}

This is neither the ``return time of a signal'' nor a physical echo. It is the
limiting time for null propagation to render every point of the spatial section
causally accessible from every other point in a static compact model. In the
causal coherence framework~\cite{Preece2026causal}, this is expressed through the
compact causal diameter $D_c = \pi R$, the causal accessibility ratio $\GammaT$,
and the condition $\Gamma(T) = 1$ for $T \geq \pi R/c$.

The fraction of causally accessible volume at $0 < T < t_{\rm coh}$ is
\begin{equation}
  \GammaT = \frac{cT/R - \tfrac{1}{2}\sin(2cT/R)}{\pi}.
\end{equation}
This function grows monotonically from $0$ to $1$. Thus, compact $S^3$ geometry
has not infinite causal openness, but finite global causal accessibility.

\section{Research Programme}
\label{sec:programme}

This calculation opens not a proof, but a research programme.

If the small negative $\Omega_K$ is accepted as a possible imprint of real positive
spatial curvature, one can construct models not in infinite $\mathbb{R}^3$ but in
compact $S^3_R$ geometry with $R = 548.3\Gly$. The following
directions are natural.


\subsection{CMB Low Multipoles and the Compact-Mode Test}
\label{sec:cmb-low-l}
The compact $S^3_R$ reconstruction naturally motivates a dedicated analysis of
the largest angular scales of the cosmic microwave background. The observed CMB
temperature field exhibits several well-known low-multipole anomalies, including
suppressed quadrupole power, quadrupole--octupole alignment, unusually weak
two-point correlations at angular separations above roughly $60^\circ$, and
related large-scale asymmetry or parity features~\cite{Copi2010,PlanckIsotropy2019}.
These anomalies are not regarded here as proof of nontrivial topology. Their
statistical significance is modest and is affected by cosmic variance, foreground
modelling, sky cuts, and look-elsewhere effects. They are therefore used only as
phenomenological targets for a compact-mode test.

The correct predictive object is the closed-FLRW hyperspherical transfer
spectrum, not the slogan that compactness suppresses the quadrupole. For
$X,Y\in\{T,E\}$ the compact calculation should take the schematic form
\begin{equation}
\boxed{
  C_\ell^{XY,S^3}
  =\sum_\beta W_\beta\,
  \mathcal P_{\mathcal R}(\beta)
  \Delta_{\ell\beta}^{X}\Delta_{\ell\beta}^{Y}.}
\end{equation}
The flat-space comparison is
\begin{equation}
  C_\ell^{XY,\mathbb R^3}
  =4\pi\int d\ln k\,
  \mathcal P_{\mathcal R}(k)
  \Delta_\ell^X(k)\Delta_\ell^Y(k).
\end{equation}
The compact-mode prediction is the difference
\begin{equation}
  \Delta C_\ell^{XY}=C_\ell^{XY,S^3}-C_\ell^{XY,\mathbb R^3}.
\end{equation}
A likelihood-level test should use the data covariance:
\begin{equation}
  \Delta\chi^2
  =\sum_{\ell\ell',XY}
  \Delta C_\ell^{XY}
  \mathrm{Cov}^{-1}_{\ell XY,\ell'X'Y'}
  \Delta C_{\ell'}^{X'Y'} .
\end{equation}
This is the operative CMB prediction pipeline.

For the fiducial radius used in this paper,
\begin{equation}
  R_{\mathcal S^3}=548.3\Gly,
\end{equation}
the observed particle-horizon radius subtends only
\begin{equation}
  \chi_{\rm obs}
  =\frac{D_{\rm particle}}{R_{\mathcal S^3}}
  =0.08412\,\mathrm{rad}\simeq4.82^\circ .
\end{equation}
Thus the observable region is a small quasi-local patch of the global carrier.
The low multipoles cannot be claimed to be suppressed by topology alone without
an explicit transfer calculation. Evidence in favour of the compact-mode
interpretation would consist not merely of low quadrupole power, but of a
coherent pattern across $TT$, $TE$, and $EE$, together with phase or correlation
structures characteristic of the selected sector. Temperature-only anomalies are
therefore insufficient; polarization is decisive~\cite{LiteBIRD2021}.

\subsection{Large-Scale Structure}

One may investigate whether a compact, gently curved geometry creates preferential
conditions for large arcs, rings, walls, and correlation structures beyond the
standard homogeneity scale. The formal chain is:
\begin{equation}
  \Omega_K \mapsto R_{\mathcal{S}^3}^{\rm global}
  \mapsto (S^3, g_R)
  \mapsto \Delta_{S^3}
  \mapsto Y_{n\ell\mu}
  \mapsto \text{projected shells}
  \mapsto \text{rings/arcs.}
\end{equation}

\subsection{Compact-Mode vs. Box-Mode Simulations}

A periodic Euclidean box $T^3$ is not equivalent to $S^3$. They have different
modes, different spectral bases, and different global geometry. Simulations on
$S^3$ harmonics can therefore produce different infrared effects.

\subsection{Matter and Gravity Dynamics}

The next step is to specify the dynamics: FLRW evolution on $S^3$, perturbation
theory, transfer functions, structure growth, lensing, and the CMB angular power
spectrum. These are dynamical tasks, not consequences of topology alone. The
compact carrier supplies the background geometry and spectral catalogue; it does
not by itself determine amplitudes, matter content, or an effective acceleration
sector.

\subsection{Causal Accessibility Observables}

Through $\GammaT$, $D_c = \pi R$, and compact null accessibility, one can search
not for ``echoes'' but for correlation signatures of finite global causal
accessibility.

\subsection{Global/Local Spectral Bridge to Ring-Like Ultra-Large-Scale Structures}
\label{sec:rings}

In~\cite{Preece2026rings} we identified ring-like and arc-like structures at
$z \approx 0.8$ with an effective spectral coherence scale
$R_{\rm eff} \approx 474.5\,\mathrm{Mpc} \approx 1.547\Gly$.
The present curvature reconstruction requires a strict separation between two
scales that must not be conflated:
\begin{equation}
  \boxed{R_{\mathcal{S}^3}^{\rm global} \neq R_{\rm eff}.}
\end{equation}

\textbf{Global carrier:}
$R_{\mathcal{S}^3}\simeq548.3\Gly$ is the curvature radius of the three-sphere
itself. It sets the Laplacian eigenvalues, the discrete mode spectrum, and the
compact causal diameter. It is a property of the manifold.

\textbf{Local spectral coherence scale:}
$R_{\rm eff} \approx 474.5\,\mathrm{Mpc}$ is an effective scale extracted from
projected-shell correlations within a narrow redshift bin. Formally,
\begin{equation}
  R_{\rm eff}=R_{\rm eff}(z,n,\Delta z,{\cal P}),
\end{equation}
where ${\cal P}$ denotes the projection and selection kernel of the observed
shell. It is a property of a particular $S^3$ mode projected onto the observable
redshift shell; it depends on the shell's comoving distance, width,
hyperspherical harmonic index, and survey window. This paper does not derive
$R_{\rm eff}$; it states the geometric separation between the global carrier
radius and a projected local coherence scale.

The logical chain connecting the two levels is:
\begin{equation}
\begin{split}
  \Omega_K
  &\mapsto R_{\mathcal{S}^3}^{\rm global}
   \mapsto (S^3, g_R)
   \mapsto \Delta_{S^3} \\
  &\mapsto Y_{n\ell\mu}
   \mapsto \text{projected modes} \\
  &\mapsto R_{\rm eff}(z,n,\Delta z,{\cal P}).
\end{split}
\end{equation}
The $474.5\,\mathrm{Mpc}$ scale is therefore not the radius of the Universe
but a \emph{derived observable} of a specific compact mode in a specific shell.
The global $S^3$ carrier provides the mode catalogue; the local $R_{\rm eff}$
is a read-out of one mode in one shell.
As a first diagnostic, Sec.~\ref{sec:transfer-likelihood-compressed} records the
projected-shell estimator $\beta_{\rm shell}\simeq618.38$ for $z=0.8$ under the
physical-scale convention stated in Sec.~\ref{sec:transfer-likelihood-compressed}.
This is a target for the full shell-transfer pipeline, not a proof of a mode assignment by itself.

Before the present paper, the $S^3$ framework explained the \emph{morphology}
of ring-like structures. The curvature reconstruction now adds an
\emph{external geometric criterion}: the global carrier is conditionally
recoverable from $\Omega_K$, not just postulated for spectral convenience. In
short:
\begin{quote}
Before: $S^3$ explained ring morphology.\par
Now: $\Omega_K$ provides the reconstructable background for that morphology.
\end{quote}

\section{Canonical Reconstruction Summary}
\label{sec:canonical}

\subsection{Three-Sphere Geometry}

\begin{align}
  \mathcal{S}^3_{\rm Light}:\quad R &= 548.3\Gly, \\
  t_R &= 548.3\Gyr, \\
  D_{\rm embed}=2R &= 1096.6\Gly, \\
  D_{\rm causal}=\pi R &= 1722.6\Gly, \\
  C=2\pi R &= 3445.2\Gly, \\
  V &\simeq 3.254\times 10^{9}\Gly^3.
\end{align}

\subsection{Manifested History}

\begin{align}
  t_0 &= 13.788\Gyr, \\
  \frac{t_0}{t_R} &= 0.02515, \\
  \theta_0 &= 1.441^\circ.
\end{align}

\subsection{Observable Horizon}

\begin{align}
  D_{\rm particle} &= 46.13\Gly, \\
  \theta_{\rm particle} &= 4.820^\circ, \\
  \Gamma_{\rm obs} &= \Gamma(\chi_{\rm particle}) \approx 0.0126\%.
\end{align}

\subsection{Fundamental Geometric Relations}

The following relations hold as exact theorems of $S^3_R$ geometry:
\begin{equation}
\begin{aligned}
  D_{\rm causal} &= \pi R, &
  V_{S^3} &= 2\pi^2 R^3, \\
  K &= \frac{1}{R^2}, &
  {}^{(3)}R &= \frac{6}{R^2}.
\end{aligned}
\end{equation}
Two-scale distinction (never conflate):
\begin{equation}
\begin{aligned}
  D_{\rm embed} &= 2R
    &&\text{(in }\mathbb{R}^4\text{)}, \\
  D_{\rm causal} &= \pi R
    &&\text{(intrinsic on }S^3\text{)}.
\end{aligned}
\end{equation}




\section{Invariant Registry, Diagnostics, and Necessity Taxonomy}
\label{sec:invariant-registry-compressed}

This section consolidates the invariant content needed for the reconstruction. It is
not an encyclopedia of natural phenomena. It is a registry of the structures actually
used by this paper.

\subsection{Exact round-\texorpdfstring{$S^3_R$}{S3R} invariants}

The round carrier is homogeneous and isotropic:
\begin{equation}
  S^3_R\simeq SO(4)/SO(3),
  \qquad \operatorname{Isom}^+(S^3_R)\simeq SO(4).
\end{equation}
The stabilizer $SO(3)$ acts transitively on the unit sphere in each tangent space.
Therefore there is no nonzero fully invariant direction:
\begin{equation}
\begin{aligned}
  \mathfrak X(S^3_R)^{SO(4)}&=0,
  \qquad
  \Omega^1(S^3_R)^{SO(4)}=0,
  \qquad
  \Omega^2(S^3_R)^{SO(4)}=0,\\
  \Omega^0(S^3_R)^{SO(4)}&=\mathbb R,
  \qquad
  \Omega^3(S^3_R)^{SO(4)}=\mathbb R\,\mathrm{vol}_{S^3}.
\end{aligned}
\label{eq:full_no_direction_law}
\end{equation}
Equivalently, invariant directions, axes, ring normals, Hopf-frame directions, and
preferred frames do not follow from the round simply connected carrier alone. They can
only enter after additional structure is chosen: an observer, field configuration,
subgroup, quotient topology, survey window, or initial condition. The same isotropy
also fixes the invariant symmetric-tensor sector:
\begin{equation}
  \mathrm{Sym}^2(T^*S^3_R)^{SO(4)}=\mathbb R\,g .
\end{equation}
Thus the round carrier alone cannot generate a geometric anisotropic stress tensor;
any anisotropy must come from perturbations, matter fields, quotient topology, survey
selection, or initial data.

The local curvature catalogue is completely fixed by $R$:
\begin{equation}
\begin{aligned}
  R_{ijkl} &= R^{-2}(g_{ik}g_{jl}-g_{il}g_{jk}),\\
  R_{ij} &= 2R^{-2}g_{ij},
  \qquad {}^{(3)}R = 6R^{-2},
  \qquad \nabla_mR_{ijkl}=0 .
\end{aligned}
\end{equation}
The scalar curvature polynomials are therefore powers of $R^{-2}$, for example
\begin{equation}
  R_{ijkl}R^{ijkl}=12R^{-4},
  \qquad
  R_{ij}R^{ij}=12R^{-4}.
\end{equation}
In three spatial dimensions the Weyl tensor has no independent components; the
conformal-flatness diagnostic is the Cotton tensor. For the round carrier,
\begin{equation}
  C^{\rm Cotton}_{ijk}=0,
\end{equation}
so the intrinsic spatial curvature is fully fixed by the metric and the single radius
$R$. This statement is purely spatial and does not assert that the four-dimensional
perturbed FLRW spacetime has vanishing Weyl curvature.

\subsection{Compact spectral diagnostics}

The empirical content of $S^3_R$ mode discreteness is the replacement of a flat
continuum by a compact hyperspherical catalogue,
\begin{equation}
  \int d^3k \quad\longrightarrow\quad \sum_{\beta\ell m},
\end{equation}
with
\begin{equation}
\begin{aligned}
  -\Delta Q_{\beta\ell m}
  &=\frac{\beta^2-1}{R^2}Q_{\beta\ell m},\\
  g^{(0)}_\beta&=\beta^2,
  \qquad
  k_\beta=\frac{\sqrt{\beta^2-1}}{R}.
\end{aligned}
\end{equation}
For the fiducial benchmark,
\begin{equation}
\begin{aligned}
  R^{-1}&=1.8237\times10^{-3}\,\mathrm{Gly}^{-1},\\
  k_{\beta=2}&=3.1588\times10^{-3}\,\mathrm{Gly}^{-1}.
\end{aligned}
\end{equation}
The constant mode $\beta=1$ is a background mode and carries no mean-zero spatial
structure. The first non-constant scalar mode is therefore $\beta=2$, with
\begin{equation}
  k_{\rm min}^{(0)}=\frac{\sqrt 3}{R}.
\end{equation}
This is an admissibility floor in the compact scalar catalogue. It is not, by itself,
an amplitude-suppression theorem: any statement about quadrupole power or large-angle
correlations still requires a primordial spectrum, transfer functions, survey window,
and likelihood comparison.
The exact scalar addition kernel is
\begin{equation}
  K_\beta(\gamma)=
  \frac{\beta}{2\pi^2R^3}\frac{\sin(\beta\gamma)}{\sin\gamma},
  \qquad \gamma=d(x,y)/R .
\end{equation}
The corresponding scalar correlation diagnostic is
\begin{equation}
  \xi_{S^3}(\gamma)=
  \sum_{\beta=1}^{\infty}P_\beta
  \frac{\beta}{2\pi^2R^3}\frac{\sin(\beta\gamma)}{\sin\gamma}.
\end{equation}
The addition kernel obeys the normalization rule
\begin{equation}
  \int_{S^3_R}K_\beta(\gamma(x,y))\,dV_y
  =\begin{cases}
  1, & \beta=1,\\
  0, & \beta>1,
  \end{cases}
\end{equation}
which separates the constant mode from all mean-zero modes. The high-mode counting is
controlled by Weyl's law,
\begin{equation}
  N(\Lambda)=\sum_{(\beta^2-1)/R^2\leq\Lambda}\beta^2
  \sim \frac{R^3}{3}\Lambda^{3/2}\qquad(\Lambda\to\infty).
\end{equation}
Consequently compact spectral effects are infrared diagnostics; at large $\beta$ the
compact sum approaches the flat continuum density. A useful angular scale for the
particle-horizon benchmark is
\begin{equation}
  \ell_{\rm IR}^{\rm diag}\sim\frac{1}{\chi_{\rm particle}}
  =\frac{R}{D_{\rm particle}}\simeq 11.9,
\end{equation}
but this is only a diagnostic window, not a substitute for a full hyperspherical
transfer calculation.

On a finite survey or causal domain $D$, the relevant quantity is the windowed response
\begin{equation}
  W_\beta[D]
  =\frac{1}{\mathrm{Vol}(D)^2}
  \int_DdV_x\int_DdV_y\,K_\beta(\gamma(x,y)).
\end{equation}
More generally, the finite observer domain induces a quasi-local mode-coupling matrix
\begin{equation}
\begin{aligned}
  M_{ab}[D]&=\int_D Q_a(x)Q_b^*(x)\,dV_x,\\
  a&=(\beta,\ell,m),\qquad b=(\beta^\prime,\ell^\prime,m^\prime).
\end{aligned}
\end{equation}
For $D=S^3_R$ this reduces to the orthonormal identity matrix. For a strict subdomain,
including the particle-horizon ball or a real survey mask, $M_{ab}[D]$ need not be
diagonal. This is a windowing and projection effect, not a new physical force. It can
mix global compact modes in quasi-local estimators, but it does not by itself force a
quadrupole deficit, an octupole alignment, or a preferred axis. Those claims require
$P_\beta$, transfer functions, foreground and mask modelling, and a likelihood test.
Thus compact-mode claims become observationally meaningful only after applying the
survey mask, transfer functions, and likelihood.

The untwisted Dirac spectrum is also discrete,
\begin{equation}
  \lambda_k^\pm=\pm\frac{k+3/2}{R},
  \qquad d_k=(k+1)(k+2),
\end{equation}
which gives the benchmark infrared spinor floor
\begin{equation}
  \Delta_D=5.7058\times10^{-35}\,\mathrm{eV}.
\end{equation}
This is a mathematical infrared threshold of the round spin carrier, not a neutrino-mass
prediction.

\subsection{Necessity taxonomy}

The paper uses three necessity classes. First, structures native to the full round
simply connected carrier include
\begin{equation}
  S^3\simeq SU(2),
  \qquad
  S^1\hookrightarrow S^3\to S^2,
  \qquad
  \pi_3(S^2)\simeq\mathbb Z,
\end{equation}
and the Skyrme degree
\begin{equation}
  U:S^3\to SU(2)\simeq S^3,
  \qquad \deg(U)\in\mathbb Z.
\end{equation}
Second, compactness-dependent but non-exclusive effects include discrete spectra,
finite-volume quantum-field-theory corrections, heat kernels, and Casimir-type
renormalized sums. These also occur on other compact manifolds or in bounded boxes.
Third, ordinary local physics---local quantum mechanics, chemistry, local QFT, local GR,
and laboratory-scale physics---does not require global $S^3$. It is preserved because
$U\subset S^3_R$ is locally $\mathbb R^3$.

\section{Transfer and Likelihood Pipeline}
\label{sec:transfer-likelihood-compressed}

The distance sector is the first observational layer. With $E(z)=H(z)/H_0$ supplied by
an external background expansion model,
\begin{equation}
  \chi(z)=\sqrt{|\Omega_K|}\int_0^z\frac{dz'}{E(z')},
  \qquad
  D_M(z)=R_c\sin\chi(z).
\end{equation}
The reconstruction test is
\begin{equation}
  \Omega_K(z)=\frac{[H(z)D_M'(z)]^2-c^2}{H_0^2D_M^2(z)}.
\end{equation}
A closed $S^3_R$ reconstruction requires $\Omega_K(z)<0$ and $d\Omega_K/dz=0$ within
errors. The BAO/SN/CMB distance observables are then
\begin{equation}
  D_A=\frac{D_M}{1+z},
  \qquad
  D_L=(1+z)D_M,
  \qquad
  \frac{dV}{dz\,d\Omega}=\frac{cD_M^2}{H(z)}.
\end{equation}

The CMB and large-scale-structure layer is transfer-grade rather than completed here:
\begin{equation}
  C_\ell^{XY,S^3}
  =\sum_\beta W_\beta\,
  \mathcal P_{\mathcal R}(\beta)
  \Delta_{\ell\beta}^{X}\Delta_{\ell\beta}^{Y},
  \qquad X,Y\in\{T,E\}.
\end{equation}
For a quotient topology $S^3/\Gamma$, the mode catalogue is projected by the deck group:
\begin{equation}
  \mathcal H_\beta(S^3/\Gamma)=\mathcal H_\beta(S^3)^\Gamma,
  \qquad
  m_\beta^\Gamma=\frac1{|\Gamma|}\sum_{g\in\Gamma}\chi_\beta(g).
\end{equation}
The full harmonic covariance is then
\begin{equation}
  C_{\ell m,\ell'm'}^{XY,\Gamma}
  =\sum_{\beta,s}
  \mathcal P_\beta\Delta_{\ell\beta}^X\Delta_{\ell'\beta}^Y
  \xi_{\beta s;\ell m}\xi^*_{\beta s;\ell'm'}.
\end{equation}
For the round simply connected sector with statistically isotropic perturbations this
reduces to the diagonal $C_\ell$ form. For quotient sectors it can generate off-diagonal
covariances and matched-circle-type constraints.

The registered distinguishability targets are likelihood and evidence differences,
including the Gaussian covariance KL diagnostic
\begin{equation}
  D_{\rm KL}(i\|0)=\frac12
  \left[\operatorname{tr}(C_0^{-1}C_i)-N
  +\ln\frac{\det C_0}{\det C_i}\right].
\end{equation}
Until the transfer functions, nuisance parameters, window operators, and covariance are
implemented, the compact spectrum remains a mode catalogue and diagnostic registry, not a
completed CMB or LSS prediction.

A useful length-to-mode diagnostic is
\begin{equation}
  \beta_{\rm wave}(L_{\rm com})=
  \sqrt{1+\left(\frac{2\pi R}{L_{\rm com}}\right)^2},
  \qquad
  \beta_{\rm half}\simeq\frac{\pi R}{L_{\rm com}}.
\end{equation}
For the ring-scale benchmark at $z\simeq0.8$, the quoted value $\beta\simeq618.38$
uses the convention that $R_{\rm eff}=474.5\,\mathrm{Mpc}$ is a physical shell scale
at that redshift and is first converted to a comoving length,
\begin{equation}
\begin{aligned}
  L_{\rm com}&=(1+z)R_{\rm eff}\simeq2.786\Gly,\\
  \beta_{\rm half}&\simeq\frac{\pi R}{L_{\rm com}}\simeq618.38 .
\end{aligned}
\end{equation}
If $474.5\,\mathrm{Mpc}$ is instead treated as an already-comoving full length, the
index would be about $1113$; if $2R_{\rm eff}$ is used as the comoving wavelength, it
would be about $557$. The convention must therefore be fixed before any shell claim is
used in a likelihood. This is a registered shell diagnostic, not a Fibonacci law or
topology proof.

\section{Guardrails and Nonclaims}
\label{sec:guardrails-compressed}

Topology and compactness supply invariant slots, spectra, and covariance structures; they
do not supply amplitudes by themselves. In particular, this manuscript does not derive or
abolish the dark sector, does not explain late-time acceleration, does not determine dark
matter, does not derive Standard Model parameters, and does not claim a preferred cosmic
axis. The background expansion $E(z)$ is an external input to the distance and transfer
pipeline.
The same rule applies to chiral gauge sectors and symmetry breaking. The Hopf fibration
makes Hopf and Skyrme topological sectors native on $S^3$, but it does not quantize the
magnetic helicity of an arbitrary primordial $U(1)$ field and does not force nonzero
$TB$ or $EB$ CMB spectra. Such parity-odd signatures would belong to a separate chiral
gauge-field model with its own action, coupling, initial state, and transfer functions.
Likewise, a vector condensate that breaks $SO(4)$ to a subgroup is an additional
dynamical hypothesis, not a consequence of the round carrier. Introducing such a field
would move the analysis beyond the minimal simply connected $S^3_R$ reconstruction
studied here.

The carrier geometry and the stress-energy state are distinct inputs: $T_{\mu\nu}$ is
defined on the spacetime geometry and does not by itself select the global topology. All
amplitude statements require a dynamics, transfer functions, covariance model, and
likelihood comparison.

\section{Status of Claims}
\label{sec:status}

\subsection{Mathematical Theorems}

The following are exact for any $S^3_R$:
\begin{equation}
\begin{aligned}
  D_c &= \pi R, &
  V_{S^3} &= 2\pi^2 R^3, \\
  K &= \frac{1}{R^2}, &
  {}^{(3)}R &= \frac{6}{R^2}.
\end{aligned}
\end{equation}

\subsection{Conditional Reconstruction}

If a small negative $\Omega_K$ is read as a real imprint of positive curvature,
then a large finite $S^3$ radius is recovered. This is conditional on the sign
and stability of $\Omega_K$.

\subsection{Phenomenological Programme}

CMB low-$\ell$ transfer spectra, BAO/AP distance diagnostics, large-scale
structure shell projections, compact harmonic simulations, causal accessibility
correlations, and quotient-topology circle or off-diagonal covariance searches.

\subsection{What Is \emph{Not} Claimed}

\begin{itemize}
  \item Global $S^3$ topology is \emph{not} proved.
  \item The finite curvature radius $R$ or intrinsic diameter $\pi R$ is \emph{not} a spatial edge or wall.
  \item The Standard Model of particle physics is \emph{not} derived.
  \item Perturbation amplitudes are \emph{not} derived.
  \item The origin of $c$ is \emph{not} proved.
  \item The $\varphi$ coordinate is \emph{not} used as a curvature measurement or prediction.
  \item Temperature-only CMB anomalies are \emph{not} claimed to prove compact topology.
  \item No tensor-to-scalar amplitude, magnetic-helicity quantization, vector condensate, $SO(4)$ symmetry-breaking field, or Standard-Model anomaly cancellation is derived from topology alone.
  \item The finite-size Casimir benchmark is \emph{not} claimed to explain dark energy.
  \item The present reconstruction does \emph{not} derive or replace the late-time acceleration sector; late-time acceleration is treated only through the background function $E(z)$.
  \item A round simply connected $S^3_R$ is \emph{not} claimed to generate a preferred spatial direction or axis.
  \item Aharonov--Bohm phases, quantum Hall Chern numbers, non-Abelian anyons, and Chern--Simons invariants are \emph{not} claimed to require a cosmological $S^3_R$ carrier.
  \item The model is predictive only through the curvature-distance law, transfer functions, and registered likelihood comparisons.
\end{itemize}


\section{Falsification Conditions}
\label{sec:falsification}

Scientific strength requires explicit falsification criteria. Each component of
this programme is falsifiable independently. For clarity, the registered targets
are separated from diagnostics:
\begin{center}
\small
\renewcommand{\arraystretch}{1.16}
\resizebox{0.98\textwidth}{!}{%
\begin{tabular}{llll}
\toprule
Target & Observable & Required calculation & Falsification condition \\
\midrule
Curvature anchor & $\Omega_K(z)$ & $H(z)$ and $D_M(z)$ reconstruction & no constant negative $\Omega_K(z)$ within errors \\
BAO geometry & $D_M/r_d$, $D_H/r_d$, $D_V/r_d$, $F_{\rm AP}$ & same priors/covariance as flat baseline & no likelihood gain or inconsistent sign under joint fit \\
Compact CMB transfer & $TT$, $TE$, $EE$ low-$\ell$ covariance & hyperspherical Boltzmann transfer sum & no $S^3$-specific improvement after parameter penalty \\
Quotient topology & matched circles or off-diagonal $C_{\ell m,\ell'm'}$ & group-projected spectrum for $S^3/\Gamma$ & positive quotient signature rejects $\Gamma=\{1\}$ minimal sector \\
Shell spectra & $\beta(z,L)$ in LSS catalogues & windowed shell projection with survey mask & no reproducible index across independent masks/catalogues \\
\bottomrule
\end{tabular}%
}
\end{center}

\subsection{Global Curvature Reconstruction}

The global closed-$S^3$ reading rests entirely on a robustly non-zero curvature
modulus together with a closed-sign interpretation, $\Omega_K^{S^3}<0$.
Falsification condition: if improved precision from DESI, Euclid, CMB-S4, or
successor surveys drives $|\Omega_K|\to0$, then $R_c\to\infty$ and the finite
radius reconstruction loses its observational anchor. If future surveys find
$|\Omega_K| < 10^{-4}$ at high significance, or if the reconstructed
$\Omega_K(z)$ of Sec.~\ref{sec:distance-law} is not redshift-independent within
errors, the specific fiducial radius $R_{\mathcal{S}^3} = 548\Gly$ loses its
observational anchor. The compact-mode phenomenology may remain as an independent
hypothesis, but the RCP linkage breaks.

\subsection{CMB Low-Multipole Prediction}

The $S^3$ Laplacian spectrum supplies a discrete compact-mode catalogue, but for
the large fiducial radius used here it does not by itself imply a hard infrared
cutoff in the observed low multipoles. The test is therefore not simply whether
the quadrupole is low. The test is whether a closed-FLRW hyperspherical transfer
calculation predicts a coherent, $S^3$-specific pattern across low-$\ell$ $TT$,
$TE$, and $EE$ spectra, phase correlations, and off-diagonal mode couplings.

Falsification condition: if improved CMB temperature and polarization data, for
example from LiteBIRD or CMB-S4, show no detectable $S^3$-specific low-$\ell$
pattern under the fiducial curvature radius, the compact-mode CMB explanation is
weakened. Temperature-only anomalies are insufficient; polarization is decisive.
A concrete likelihood implementation should report
\begin{equation}
  \Delta\chi^2_{\rm CMB}
  =\chi^2_{S^3\text{-transfer}}-\chi^2_{\Lambda{\rm CDM}}
\end{equation}
for the registered low-$\ell$ $TT/TE/EE$ data vector. If the compact transfer
model fails to improve the fit, or is disfavoured after the appropriate
information-criterion penalty, the CMB compact-mode claim fails as a predictive
explanation.

A positive detection of matched circles with a quotient-specific phase relation
would favour a multiply connected spherical space form such as $S^3/\Gamma$ and
would falsify the minimal simply connected sector choice, although not the broader
positive-curvature compact-geometry programme.

\subsection{Ring-Like Structure Predictions}

The ring analysis of~\cite{Preece2026rings} predicts specific angular scales,
redshift-dependent $R_{\rm eff}(z,n,\Delta z,{\cal P})$, and a discrete mode
selection rule. Falsification conditions:
\begin{itemize}
  \item DESI/Euclid find no preferred angular clustering at the predicted
    compact-mode scales in the relevant redshift bins
    $\Rightarrow$ ring-mode assignment weakened.
  \item The observed ring scales at $z \approx 0.8$ are better explained by
    acoustic or projection effects within $\Lambda$CDM
    $\Rightarrow$ $S^3$-mode interpretation loses parsimony.
  \item A claimed shell index $\beta(z,L)$ is not reproduced in independent
    catalogues, redshift shells, or survey masks
    $\Rightarrow$ projected spectral assignment falsified.
\end{itemize}

\subsection{Compact vs. Box Simulations}

Falsification condition: if $S^3$-harmonic N-body simulations produce no
infrared differences from $T^3$ periodic-box simulations at the observed scales,
the compact-geometry simulation programme is weakened.

\medskip
\noindent\textbf{Summary of falsification logic:} Each of the five components
is independently testable. Falsification of the global reconstruction
($\Omega_K \to 0$) does not automatically falsify the local spectral programme
(ring modes, compact simulations). The programme is \emph{modular}: observational
pressure on any one component produces knowledge, not collapse.


\paragraph*{Carrier-state distinction.}
The carrier geometry and the stress-energy state are distinct inputs: $T_{\mu\nu}$ is
specified on the spacetime geometry and does not by itself select the global topology.

\section{Conclusion}
\label{sec:conclusion}

If a small curvature modulus is used together with the closed-sign interpretation
$\Omega_K^{S^3}<0$, it specifies a vast but finite $S^3$ geometry with curvature
radius $R \approx 548.3\Gly$.

In such a geometry, the observable Universe occupies only a small quasi-local
region of approximately $4.82^\circ$ in particle-horizon radius along the great
geodesic of the three-sphere. The locally nearly flat reconstruction is therefore
natural, but it does not prove a global $\mathbb{R}^3$ ontology.

The principal scientific claim is:
\begin{quote}
The flat infinite framework is not mandatory. The compact $S^3$ framework
preserves local physics while yielding testable global consequences.
\end{quote}

This makes the reconstruction scientifically substantive: it does not ``invent an
age for the Universe'', but recovers a possible curvature scale of a compact
three-sphere and thereby provides a numerical foundation for further CMB modelling,
large-scale structure analysis, and compact spectral cosmology. The CMB
low-multipole anomalies do not prove this geometry; they define a disciplined
compact-mode falsification target for future temperature and polarization data.

\clearpage
\section*{Summary of Key Numerical Values}

\begin{center}
\textbf{Table I. Fiducial benchmark numerical values of the $S^3_R$ reconstruction
($R = 548.324\,513\,026\,856\,\mathrm{Gly}$). These machine-precision values are
algebraic outputs of the chosen inputs, not observational significant figures.
Distances are in gigalight-years (Gly), times in gigayears (Gyr).}
\label{tab:summary}
\small
\renewcommand{\arraystretch}{1.18}
\resizebox{0.98\textwidth}{!}{%
\begin{tabular}{llll}
\toprule
Quantity & Symbol & Value & Unit \\
\midrule
Curvature radius & $R_{\mathcal{S}^3}$ & $548.324\,513\,026\,856$ & Gly \\
Light-travel time of $R$ & $t_{\mathcal{S}^3}$ & $548.324\,513\,026\,856$ & Gyr \\
Embedding diameter & $D_{\rm embed}=2R$ & $1096.649\,026\,053\,712$ & Gly \\
Causal intrinsic diameter & $D_{\rm causal}=\pi R$ & $1722.612\,261\,908\,368$ & Gly \\
Full great circle & $C=2\pi R$ & $3445.224\,523\,816\,736$ & Gly \\
Proper volume & $V_{S^3}=2\pi^2R^3$ & $3.254\,188\,648\,717\times10^{9}$ & Gly$^3$ \\
Sectional curvature & $K=1/R^2$ & $3.326\,018\,624\,722\times10^{-6}$ & Gly$^{-2}$ \\
Scalar curvature & ${}^{(3)}R=6/R^2$ & $1.995\,611\,174\,833\times10^{-5}$ & Gly$^{-2}$ \\
Causal accessibility time & $t_{\rm coh}=\pi R/c$ & $1722.612\,261\,908\,368$ & Gyr \\
Cosmic age & $t_0$ & $13.787\,718\,785\,602$ & Gyr \\
Age ratio & $t_0/t_{\mathcal{S}^3}$ & $0.025\,145\,180\,377\,75$ & $2.515\%$ \\
Age as great-circle arc & $\theta_0$ & $1.440\,712\,710\,740$ & deg \\
Particle horizon & $D_{\rm particle}$ & $46.125\,169\,765\,204$ & Gly \\
Horizon arc & $\theta_{\rm particle}$ & $4.819\,732\,647\,519$ & deg \\
Volumetric fraction & $\Gamma_{\rm obs}$ & $\approx 0.0126$ & $\%$ \\
Joint accessibility diagnostic & $\Gamma_{\rm ST}$ & $4.255\,135\times10^{-10}$ & -- \\
Boundary-area benchmark & $A_{\partial D_{\rm obs}}$ & $2.667\,234\times10^{4}$ & Gly$^2$ \\
Dirac infrared gap & $\Delta_D$ & $5.705\,785\times10^{-35}$ & eV \\
Conformal scalar Casimir density & $\rho_{\rm Cas}^{(0)}$ & $9.215\,238\times10^{-141}$ & J m$^{-3}$ \\
Projected shell estimator & $\beta_{\rm shell}(z=0.8)$ & $618.38$ & convention-dependent \\
Full-wavelength shell bridge & $\beta_{\rm wave}(L_{\rm com})$ & $\sqrt{1+(2\pi R/L_{\rm com})^2}$ & -- \\
Quasi-local scalar kernel & $\xi_{S^3}(\gamma)$ & discrete $\beta$ sum & -- \\
\bottomrule
\end{tabular}%
}
\end{center}



\clearpage

\appendix
\section{Derivation of the Scalar Addition Theorem on \texorpdfstring{$S^3_R$}{S3R}}
\label{app:addition-theorem}

This appendix derives the scalar addition theorem used in the compact spectral
kernel. Let $\mathcal H_\beta$ denote the eigenspace of the scalar
Laplace--Beltrami operator on the round three-sphere $S^3_R$ with eigenvalue
\begin{equation}
  -\Delta_{S^3}Q_{\beta\ell m}
  =
  \frac{\beta^2-1}{R^2}Q_{\beta\ell m},
  \qquad
  \beta=1,2,\ldots,
\end{equation}
where $\ell=0,\ldots,\beta-1$ and $m=-\ell,\ldots,\ell$. The dimension of this
eigenspace is
\begin{equation}
  \dim\mathcal H_\beta=\beta^2 .
\end{equation}
For any orthonormal basis $\{Q_{\beta\ell m}\}$ of $\mathcal H_\beta$, define the
finite-rank spectral projector kernel
\begin{equation}
  K_\beta(x,y)
  =
  \sum_{\ell=0}^{\beta-1}
  \sum_{m=-\ell}^{\ell}
  Q_{\beta\ell m}(x)\,
  Q^*_{\beta\ell m}(y).
\end{equation}
It satisfies the reproducing property
\begin{equation}
  \int_{S^3_R}K_\beta(x,z)f(z)\,dV_z=f(x),
  \qquad f\in\mathcal H_\beta.
\end{equation}

By $SO(4)$ invariance of the round metric, the kernel can depend on $x$ and $y$
only through their geodesic separation
\begin{equation}
  \gamma=\frac{d(x,y)}{R},
  \qquad
  \cos\gamma=\frac{x\cdot y}{R^2}.
\end{equation}
Thus $K_\beta(x,y)=K_\beta(\gamma)$ is a zonal harmonic. On $S^3$, the zonal
harmonics of degree $n$ are proportional to Gegenbauer polynomials
$C_n^{(1)}(\cos\gamma)$. Since
\begin{equation}
  C_n^{(1)}(\cos\gamma)
  =
  \frac{\sin[(n+1)\gamma]}{\sin\gamma},
\end{equation}
and the convention used in this manuscript is $\beta=n+1$, the kernel must have
the form
\begin{equation}
  K_\beta(\gamma)
  =
  A_\beta
  \frac{\sin(\beta\gamma)}{\sin\gamma}.
\end{equation}

The constant $A_\beta$ is fixed by evaluating the diagonal. Since $K_\beta(x,x)$
is constant by homogeneity,
\begin{equation}
  \int_{S^3_R}K_\beta(x,x)\,dV_x
  =
  \dim\mathcal H_\beta
  =
  \beta^2 .
\end{equation}
Using $\operatorname{Vol}(S^3_R)=2\pi^2R^3$, this gives
\begin{equation}
  K_\beta(x,x)=\frac{\beta^2}{2\pi^2R^3}.
\end{equation}
Moreover,
\begin{equation}
  \lim_{\gamma\to0}
  \frac{\sin(\beta\gamma)}{\sin\gamma}
  =
  \beta,
\end{equation}
so
\begin{equation}
  A_\beta=\frac{\beta}{2\pi^2R^3}.
\end{equation}
Therefore the scalar addition theorem on the round three-sphere is
\begin{equation}
  \boxed{
  \sum_{\ell=0}^{\beta-1}
  \sum_{m=-\ell}^{\ell}
  Q_{\beta\ell m}(x)Q^*_{\beta\ell m}(y)
  =
  \frac{\beta}{2\pi^2R^3}
  \frac{\sin(\beta\gamma)}{\sin\gamma}
  } .
\end{equation}

Equivalently,
\begin{equation}
  -\Delta_x K_\beta(x,y)
  =
  \frac{\beta^2-1}{R^2}K_\beta(x,y),
\end{equation}
or
\begin{equation}
  \left(-\Delta_x-\frac{\beta^2-1}{R^2}\right)K_\beta(x,y)=0.
\end{equation}
This equation holds for the smooth finite-rank projector $K_\beta$. It should not
be confused with a Green function: a delta source appears only after summing over
all eigenspaces,
\begin{equation}
  \sum_{\beta=1}^{\infty}K_\beta(x,y)
  =
  \delta_{S^3_R}(x,y)
\end{equation}
in the distributional sense.

For the flat-space limit, set $r=R\gamma$ and $k=\beta/R$. At fixed physical
separation $r$ and $R\to\infty$,
\begin{equation}
  K_\beta(r)
  \sim
  \frac{1}{R}
  \frac{k^2}{2\pi^2}
  \frac{\sin(kr)}{kr}.
\end{equation}
The factor $1/R$ is the discrete spacing $\Delta k$. Hence the sum over $\beta$
becomes the ordinary flat radial Fourier measure,
\begin{equation}
  \sum_\beta K_\beta(r)
  \longrightarrow
  \int_0^\infty
  \frac{k^2\,dk}{2\pi^2}
  \frac{\sin(kr)}{kr}.
\end{equation}

This theorem is the exact bridge between the compact spectral catalogue and
isotropic two-point statistics on $S^3_R$. It collapses the sum over $(\ell,m)$
into a single zonal kernel indexed by $\beta$, and it underlies the compact
correlation functions, shell projections, windowed responses, mode-coupling
matrices, and transfer-function sums used in the main text.

\section{Complete Invariant Registry}
\label{app:complete-registry}

This appendix summarizes the invariant content used in the paper. It is not an
encyclopedia of all natural phenomena. It is a closed registry only with respect
to the geometric, topological, spectral, observational, and diagnostic structures
actually used in the reconstruction.

\begin{center}
\small
\renewcommand{\arraystretch}{1.15}
\resizebox{0.98\textwidth}{!}{%
\begin{tabular}{llll}
\toprule
Class & Invariant object & Formula or criterion & Role in the paper \\
\midrule
Local curvature & Sectional curvature & $K=R^{-2}$ & Exact geometry \\
Local curvature & Ricci tensor & $R_{ij}=2R^{-2}g_{ij}$ & Exact geometry \\
Local curvature & Scalar curvature & ${}^{(3)}R=6R^{-2}$ & Exact geometry \\
Local curvature & Riemann tensor & $R_{ijkl}=R^{-2}(g_{ik}g_{jl}-g_{il}g_{jk})$ & Exact geometry \\
Local curvature & Cotton tensor & $C^{\rm Cotton}_{ijk}=0$ & Spatial conformal flatness of round carrier \\
Local symmetry & Covariant curvature derivative & $\nabla_mR_{ijkl}=0$ & Round local symmetry \\
No-axis theorem & Invariant tangent vector & $(T_pS^3_R)^{SO(3)}=\{0\}$ & Excludes geometric preferred axes \\
Invariant tensors & Symmetric 2-tensor sector & $\mathrm{Sym}^2(T^*S^3_R)^{SO(4)}=\mathbb Rg$ & Excludes geometry-only anisotropic stress \\
Global geometry & Proper volume & $V_{S^3}=2\pi^2R^3$ & Compact volume \\
Global geometry & Intrinsic diameter & $D_c=\pi R$ & Maximal geodesic distance \\
Global geometry & Great circle & $C=2\pi R$ & Closed geodesic scale \\
Global geometry & Causal-volume identity & $D_cV_{S^3}^{1/3}/R^2=\pi(2\pi^2)^{1/3}$ & Pure round-$S^3$ benchmark \\
Geodesic ball & Volume fraction & $\Gamma(\chi)=[\chi-\tfrac12\sin(2\chi)]/\pi$ & Observable-domain fraction \\
Boundary benchmark & Geodesic-sphere area & $A(\chi)=4\pi R^2\sin^2\chi$ & Isoperimetric diagnostic \\
Scalar spectrum & Laplace eigenvalues & $\lambda_\beta=(\beta^2-1)/R^2$ & Compact mode catalogue \\
Scalar spectrum & Degeneracy & $g^{(0)}_\beta=\beta^2$ & Mode counting \\
Scalar spectrum & Weyl counting law & $N(\Lambda)\sim R^3\Lambda^{3/2}/3$ & High-mode flat-limit check \\
Wave number & Discrete scalar wave number & $k_\beta=\sqrt{\beta^2-1}/R$ & Bridge to transfer variables \\
Addition theorem & Scalar two-point kernel & $K_\beta(\gamma)=\dfrac{\beta}{2\pi^2R^3}\dfrac{\sin(\beta\gamma)}{\sin\gamma}$ & Exact correlation kernel \\
Addition theorem & Kernel normalization & $\int_{S^3}K_\beta\,dV=\delta_{\beta1}$ & Constant-mode/mean-zero separation \\
Spinors & Untwisted Dirac spectrum & $\lambda_k^\pm=\pm(k+\tfrac32)/R$ & Infrared spinor gap \\
Spinors & Dirac degeneracy & $d_k=(k+1)(k+2)$ & Spinor mode counting \\
Topology & Fundamental group & $\pi_1(S^3)=0$ & Simply connected sector \\
Topology & Cohomology & $H^1=H^2=0,\ H^3\simeq\mathbb Z$ & No native one/two cohomology; volume class \\
Topology & Hopf sector & $S^1\hookrightarrow S^3\to S^2$, $\pi_3(S^2)\simeq\mathbb Z$ & Native Hopf charge domain \\
Topology & Skyrme degree & $\deg(S^3\to SU(2)\simeq S^3)\in\mathbb Z$ & Native Skyrmion topology \\
Quotients & Deck-group projection & $\mathcal H_\beta(S^3/\Gamma)=\mathcal H_\beta(S^3)^\Gamma$ & Alternative topology sector \\
Quotients & Mode multiplicity & $m_\beta^\Gamma=|\Gamma|^{-1}\sum_{g\in\Gamma}\chi_\beta(g)$ & Quotient spectral selection \\
FLRW distance & Transverse distance & $D_M(z)=R_c\sin\chi(z)$ & Observational curvature law \\
FLRW test & Curvature consistency & $\Omega_K(z)=([H D_M']^2-c^2)/(H_0^2D_M^2)$ & Reconstruction test \\
Causal geometry & Accessibility time & $t_{\rm coh}=\pi R/c$ & Global null-accessibility scale \\
Causal geometry & Accessibility fraction & $\Gamma(T)=[cT/R-\tfrac12\sin(2cT/R)]/\pi$ & Causal-domain fraction \\
Spectral bridge & Full-wave convention & $\beta_{\rm wave}=\sqrt{1+(2\pi R/L_{\rm com})^2}$ & Diagnostic length-to-mode map \\
Spectral bridge & Half-wave convention & $\beta_{\rm half}\simeq \pi R/L_{\rm com}$ & Registered shell diagnostic \\
Windowed response & Domain-mode response & $W_\beta[D]=\mathrm{Vol}(D)^{-2}\int_DdV_x\int_DdV_yK_\beta(\gamma)$ & Mask/domain diagnostic \\
Likelihood chain & Predictive endpoint & $\Delta\ln\mathcal L,\ \Delta\ln Z$ & Required status for amplitude predictions \\
Guardrail & Non-prediction criterion & no likelihood slot $\Rightarrow$ diagnostic only & Prevents overclaiming \\
\bottomrule
\end{tabular}%
}
\end{center}

The registry separates exact mathematical statements from observational claims.
Objects in the first columns are exact within the chosen round-$S^3_R$ model.
Their physical use becomes predictive only after they enter an observable transfer
or likelihood pipeline. The causal-volume identity and the shell-index rows are
benchmarks: they are not independent detections of topology. This is the operative
closure rule of the paper:
\begin{equation}
\begin{aligned}
  S^3_R
  &\longrightarrow
  \{D_M,H,dV/dz\,d\Omega,k_\beta,K_\beta\} \\
  &\longrightarrow
  \{C_\ell^{TT,TE,EE},P(k,z),\xi(r,z)\} \\
  &\longrightarrow
  \Delta\ln\mathcal L\ \text{or}\ \Delta\ln Z .
\end{aligned}
\end{equation}
Anything outside this chain is retained only as a diagnostic, a taxonomy entry,
or an explicit non-claim.

\clearpage

\section{Optional Descriptive \texorpdfstring{$\varphi$}{phi} Scale Coordinate}
\label{app:phi-coordinate}

The $\varphi$-hierarchy is not used as evidence for the compact reconstruction. It is a
descriptive scale coordinate only:
\begin{equation}
  n_\varphi(R)=\frac{\ln(R/\ell_P)}{\ln\varphi},
  \qquad
  \varphi=\frac{1+\sqrt5}{2}.
\end{equation}
For the fiducial radius this gives $n_\varphi(R_{\mathcal S^3})\simeq299.09$. This
number is not a topology proof, not a curvature measurement, and not an amplitude
prediction. The established appearance of $\varphi$ in Fibonacci anyon theory is a
mathematical analogy in topological quantum matter, not an independent derivation of
cosmological curvature.

\begin{thebibliography}{99}

\bibitem{Preece2026causal}
A.~Preece and B.~Batenin,
``Causal Coherence and Compact Null Accessibility in
  $S^3_R \times \mathbb{R}$ Cosmology,''
Independent Researchers, Berlin (2026), manuscript in preparation.

\bibitem{Preece2026rings}
A.~Preece and B.~Batenin,
``Ring-Like Ultra-Large-Scale Structures at $z \approx 0.8$:
  A Compact $S^3$-Mode Analysis,''
Independent Researchers, Berlin (2026), manuscript in preparation.

\bibitem{Preece2026suspension}
A.~Preece and B.~Batenin,
``Suspension Hierarchy $S^0 \to S^1 \to S^2$
  and $S^3 \to S^4$ Self-Similar Scale Organisation,''
Independent Researchers, Berlin (2026), manuscript in preparation.

\bibitem{Planck2020}
N.~Aghanim \textit{et al.} (Planck Collaboration),
``Planck 2018 results. VI. Cosmological parameters,''
Astron.\ Astrophys.\ \textbf{641}, A6 (2020),
doi:10.1051/0004-6361/201833910.

\bibitem{PlanckIsotropy2019}
Y.~Akrami \textit{et al.} (Planck Collaboration),
``Planck 2018 results. VII. Isotropy and statistics of the CMB,''
Astron.\ Astrophys.\ \textbf{641}, A7 (2020),
arXiv:1906.02552.


\bibitem{ACTDR6Extensions}
E.~Calabrese \textit{et al.} (ACT Collaboration),
``The Atacama Cosmology Telescope: DR6 Constraints on Extended Cosmological Models,''
arXiv:2503.14454 (2025).

\bibitem{PlanckTopology2016}
P.~A.~R.~Ade \textit{et al.} (Planck Collaboration),
``Planck 2015 results. XVIII. Background geometry and topology of the Universe,''
Astron.\ Astrophys.\ \textbf{594}, A18 (2016),
arXiv:1502.01593.

\bibitem{COMPACTPromise2024}
N.~J.~Cornish \textit{et al.} (COMPACT Collaboration),
``Promise of Future Searches for Cosmic Topology,''
arXiv:2210.11426 (2024 version).

\bibitem{COMPACTPartIIA2024}
J.~R.~Eskilt \textit{et al.} (COMPACT Collaboration),
``Cosmic topology. Part IIa. Eigenmodes, correlation matrices, and detectability
  of orientable Euclidean manifolds,''
JCAP \textbf{03}, 036 (2024), arXiv:2306.17112.


\bibitem{COMPACTPower2026}
J.~Carron Duque \textit{et al.} (COMPACT Collaboration),
``Cosmic topology. Part IIc. Detectability with non-standard primordial power spectrum,''
arXiv:2602.15527 (2026).

\bibitem{COMPACTLens2025}
S.~Saha \textit{et al.} (COMPACT Collaboration),
``Cosmic topology. Part Ic. Limits on lens spaces from circle searches,''
JCAP \textbf{01}, 004 (2025), arXiv:2409.02226.

\bibitem{DESI2024}
A.~G.~Adame \textit{et al.} (DESI Collaboration),
``DESI 2024 VI: Cosmological Constraints from the Measurements of
  Baryon Acoustic Oscillations,''
arXiv:2404.03002 (2024).

\bibitem{DESIDR2}
M.~Abdul-Karim \textit{et al.} (DESI Collaboration),
``DESI DR2 Results II: Measurements of Baryon Acoustic Oscillations and
  Cosmological Constraints,''
arXiv:2503.14738 (2025).

\bibitem{Euclid2024}
A.~Blanchard \textit{et al.} (Euclid Collaboration),
``Euclid preparation: Sensitivity to dark energy,''
Astron.\ Astrophys.\ \textbf{690}, A148 (2024).

\bibitem{LiteBIRD2021}
E.~Allys \textit{et al.} (LiteBIRD Collaboration),
``Probing cosmic inflation with the LiteBIRD cosmic microwave background polarization survey,''
Prog.\ Theor.\ Exp.\ Phys.\ \textbf{2023}, 042F01 (2023),
arXiv:2202.02773; see also arXiv:2101.12449 for the mission overview.

\bibitem{Luminet2003}
J.-P.~Luminet, J.~R.~Weeks, A.~Riazuelo, R.~Lehoucq, and J.-P.~Uzan,
``Dodecahedral space topology as an explanation for weak wide-angle temperature
  correlations in the cosmic microwave background,''
Nature \textbf{425}, 593--595 (2003).

\bibitem{Cornish1998}
N.~J.~Cornish, D.~N.~Spergel, and G.~D.~Starkman,
``Measuring the topology of the Universe,''
Proc.\ Natl.\ Acad.\ Sci.\ \textbf{95}, 82--84 (1998).

\bibitem{Aurich2005}
R.~Aurich, S.~Lustig, and F.~Steiner,
``CMB anisotropy of the Poincar\'e dodecahedron,''
Class.\ Quantum Grav.\ \textbf{22}, 2061--2083 (2005).

\bibitem{LachiezeReyLuminet1995}
M.~Lachi\`eze-Rey and J.-P.~Luminet,
``Cosmic topology,''
Phys.\ Rep.\ \textbf{254}, 135--214 (1995).

\bibitem{Weeks2003}
J.~R.~Weeks,
``Detecting topology in a nearly flat hyperbolic universe,''
Mod.\ Phys.\ Lett.\ A \textbf{18}, 2099--2108 (2003).

\bibitem{Handley2021}
W.~Handley,
``Curvature tension: Evidence for a closed universe,''
Phys.\ Rev.\ D \textbf{103}, L041301 (2021).

\bibitem{DiValentino2020}
E.~Di~Valentino, A.~Melchiorri, and J.~Silk,
``Planck evidence for a closed Universe and a possible crisis for cosmology,''
Nature Astron.\ \textbf{4}, 196--203 (2020).

\bibitem{Copi2010}
C.~J.~Copi, D.~Huterer, D.~J.~Schwarz, and G.~D.~Starkman,
``Large-angle anomalies in the CMB,''
Adv.\ Astron.\ \textbf{2010}, 847541 (2010).

\bibitem{Bielewicz2011}
P.~Bielewicz and A.~J.~Banday,
``Constraints on the topology of the Universe derived from the
  lensing power spectrum,''
Mon.\ Not.\ R.\ Astron.\ Soc.\ \textbf{412}, 2104--2110 (2011).

\bibitem{MorganTian2007}
J.~Morgan and G.~Tian,
\textit{Ricci Flow and the Poincar\'e Conjecture},
Clay Mathematics Monographs, Vol.~3, American Mathematical Society / Clay Mathematics Institute (2007).

\bibitem{ClayPoincare}
Clay Mathematics Institute,
``Poincar\'e Conjecture,'' Millennium Problems overview,
https://www.claymath.org/millennium/poincare-conjecture/ .

\bibitem{Adams1962}
J.~F.~Adams,
``Vector fields on spheres,''
Annals of Mathematics \textbf{75}, 603--632 (1962).

\bibitem{BenedettiLisca2018}
R.~Benedetti and P.~Lisca,
``Framing 3-manifolds with bare hands,''
arXiv:1806.04991.

\bibitem{SU2LieNotes}
C.~Scaduto,
``Lie Theory as Illustrated by $SU(2)$,'' lecture notes, Stony Brook University.

\bibitem{Nayak2008}
C.~Nayak, S.~H. Simon, A.~Stern, M.~Freedman, and S.~Das Sarma,
``Non-Abelian anyons and topological quantum computation,''
Rev.\ Mod.\ Phys.\ \textbf{80}, 1083 (2008).

\bibitem{Trebst2008}
S.~Trebst, M.~Troyer, Z.~Wang, and A.~W.~W. Ludwig,
``A short introduction to Fibonacci anyon models,''
Prog.\ Theor.\ Phys.\ Suppl.\ \textbf{176}, 384--407 (2008).

\bibitem{MosseriDandoloff2001}
R.~Mosseri and R.~Dandoloff,
``Geometry of entangled states, Bloch spheres and Hopf fibrations,''
J. Phys. A: Math. Gen.\ \textbf{34}, 10243--10252 (2001),
arXiv:quant-ph/0108137.

\bibitem{Urbantke2003}
H.~K.~Urbantke,
``The Hopf fibration---seven times in physics,''
J. Geom. Phys.\ \textbf{46}, 125--150 (2003).

\bibitem{Sutcliffe2007}
P.~Sutcliffe,
``Knots in the Skyrme--Faddeev model,''
Proc. R. Soc. A \textbf{463}, 3001--3020 (2007), arXiv:0705.1468.

\bibitem{Manton1998}
N.~S.~Manton,
``Understanding Skyrmions using rational maps,''
in \textit{Proceedings of the Third European Congress of Mathematics} (1998).

\bibitem{CamporesiHiguchi1996}
R.~Camporesi and A.~Higuchi,
``On the eigenfunctions of the Dirac operator on spheres and real hyperbolic spaces,''
J. Geom. Phys.\ \textbf{20}, 1--18 (1996), arXiv:gr-qc/9505009.

\bibitem{GromovLawson1983}
M.~Gromov and H.~B.~Lawson,
``Positive scalar curvature and the Dirac operator on complete Riemannian manifolds,''
Publ. Math. IHES \textbf{58}, 83--196 (1983).


\bibitem{Witten1989}
E.~Witten,
``Quantum field theory and the Jones polynomial,''
Commun. Math. Phys. \textbf{121}, 351--399 (1989).

\bibitem{AharonovBohm1959}
Y.~Aharonov and D.~Bohm,
``Significance of electromagnetic potentials in the quantum theory,''
Phys. Rev. \textbf{115}, 485--491 (1959).

\bibitem{TKNN1982}
D.~J.~Thouless, M.~Kohmoto, M.~P.~Nightingale, and M.~den Nijs,
``Quantized Hall conductance in a two-dimensional periodic potential,''
Phys. Rev. Lett. \textbf{49}, 405--408 (1982).

\bibitem{IyerWald1994}
V.~Iyer and R.~M.~Wald,
``Some properties of Noether charge and a proposal for dynamical black hole entropy,''
Phys. Rev. D \textbf{50}, 846--864 (1994).

\bibitem{BrownYork1993}
J.~D.~Brown and J.~W.~York,
``Quasilocal energy and conserved charges derived from the gravitational action,''
Phys. Rev. D \textbf{47}, 1407--1419 (1993).

\bibitem{Mota2022}
H.~F.~S.~Mota, E.~R.~Bezerra de Mello, and K.~Bakke,
``Thermal Casimir Effect in the Einstein Universe with a Spherical Boundary,''
Universe \textbf{8}, 597 (2022).





\end{thebibliography}

\end{document}
